% Ensure theorem environments used in this chapter are available.
\makeatletter
\@ifundefined{conclusion}{\newenvironment{conclusion}[1][]{\csname remark\endcsname[#1]}{\csname endremark\endcsname}}{}
\makeatother

\chapter{From dg-Shifted Yangians to Spectral Factorization Quantum Groups}
\label{chap:dg-shifted-factorization}

\section{Orientation and first principles}

Theorem~\ref{thm:yangian-recognition} identifies the open-colour
line-side algebra $\cA^!_{\mathrm{line}}$ on the chirally Koszul locus
as a dg-shifted Yangian package whose axioms are the
$(\SCchtop)^{!}$-operadic identities (closed $=$ Lie, open $=$ Ass)
decomposed by colour. The spectral Drinfeld problem for this package is
a filtered strictification problem for the regularized holonomy of a
bar-horizontal spectral connection, not an identification of dg-shifted
Yangians with factorization quantum groups.

Three languages converge. The physical identifications of
Dimofte--Niu--Py~\cite{DNP25} are the naming of operadic data in
physical language (Remark~\ref{rem:yangian-logical-status}). The
spectral language of factorisation quantum groups is that of
Latyntsev~\cite{Latyntsev23}. The geometric language of shifted quantum
groups and critical stable envelopes is that of the shifted-Yangian
literature and the critical-stable-envelope
programme~\cite{COZZ26}.

The three languages share a single spine. A dg-shifted Yangian, viewed
through Theorem~\ref{thm:yangian-recognition} on the chirally Koszul
locus, is not \emph{already} a strict factorization quantum group: the
$\SCchtop$-algebra structure, once expressed through a bar-horizontal
strictification of its open-color projection, canonically produces a
\emph{spectral quasi-factorization quantum group}, and the passage from
quasi to strict is governed by a concrete obstruction complex whose
first low-filtration sectors are computed exactly.

The rigidity of the computed sectors is \emph{root multiplicity one}:
every root of a simple Lie algebra has multiplicity one, so each
nonzero root-sector obstruction is a scalar. This identifies the
exact test: the product-side scalar equals the coproduct-side scalar
modulo the scalar spectral coboundary. Path-type sectors satisfy this
test by the $n$-gon rigidity calculation. In general finite type, root
multiplicity one reduces strictification to the scalar obstruction map
defined below.

A strict $R$-matrix on the $A_\infty$ object would conflate colors.
The $R$-matrix is a closed-color datum (from $\FM_2(\C)$); the
$A_\infty$ structure is the open-color $E_1$-projection of the
$\SCchtop$-algebra on the chirally Koszul locus
(Theorem~\ref{thm:yangian-recognition}). The construction passes to
ordered configuration spaces (the open-color stage), produces a flat
connection, reads off regularized holonomy, and then formulates the
spectral factorization problem. In the Steinberg presentation
(Theorem~\ref{thm:steinberg-presentation}), the flat connection is the
spectral connection $\nabla_z$ of (iii)(b); strictification is the
passage from its monodromy $R(z)$ as a homotopy-coherent correspondence
to a strict factorization $R$-matrix.

\section{Complete filtered spectral setup}

\begin{definition}[Residue-bounded complete dg-shifted Yangian]
Let $Y$ be a complete separated filtered $A_\infty$-algebra
\[
Y=F^1Y\supset F^2Y\supset \cdots,
\qquad
Y\cong \varprojlim_N Y/F^{N+1}Y,
\]
equipped with:
\begin{enumerate}
\item a translation family $\tau_z$,
\item a twisted coproduct $\Delta_z$,
\item an odd Maurer--Cartan kernel
\[
r(z)=\sum_{n\ge 1} r_n z^{-n} \in (Y\widehat\otimes Y)((z^{-1})),
\]
\end{enumerate}
such that
\[
m_k(F^{p_1}Y,\dots,F^{p_k}Y)
\subset F^{p_1+\cdots+p_k+k-2}Y,
\]
and such that the residues become more convergent with pole order in the
sense that
\[
r_n(F^pY\widehat\otimes F^qY)
\subset F^{p+q+n}(Y\widehat\otimes Y).
\]
We call such an object a \emph{residue-bounded complete dg-shifted
Yangian}.
\end{definition}

The suspension of $r(z)$ to the bar convolution algebra will be denoted
by $\widetilde r(z)$; this is the degree-zero spectral kernel.

\begin{definition}[Bar-horizontal strictification]
Let $B(Y)=T^c(s^{-1}Y)$ be the completed reduced bar coalgebra of $Y$ with
bar differential $b$. A \emph{bar-horizontal strictification} of $Y$
is a complete filtered cochain complex $(V,d_V)$ quasi-isomorphic to
$B(Y)$ together with continuous degree-zero endomorphisms
\[
\Omega_{ij}(u)\in \operatorname{End}(V^{\widehat\otimes n})((u^{-1}))
\qquad (1\le i<j\le n)
\]
such that:
\begin{align*}
[d_V,\Omega_{ij}(u)]&=0,\\
[\Omega_{ij}(u),\Omega_{kl}(v)]&=0
\qquad \text{if }\{i,j\}\cap\{k,l\}=\varnothing,
\end{align*}
and for every triple $i<j<k$,
\[
[\Omega_{ij}(u),\Omega_{ik}(u+v)]
+[\Omega_{ij}(u),\Omega_{jk}(v)]
+[\Omega_{ik}(u+v),\Omega_{jk}(v)] = 0.
\]
\end{definition}

This is the additive spectral precursor of a strict spectral
$R$-matrix. It is the correct object to strictify first.

\begin{remark}[Existence of bar-horizontal strictifications]
\label{rem:bar-horizontal-existence}
For HT theories in the scope of Theorem~\ref{thm:physics-bridge}, the
perturbative expansion of the holonomy supplies candidate operators
$\Omega_{ij}(u)$ with tree-level exchange amplitude as the leading
term. Chain closure together with the disjointness and spectral Kohno
identities above produces a bar-horizontal strictification.
Theorem~\ref{thm:derived-additive-kz} derives flatness from those
identities; their existence is the input.
\end{remark}

\begin{theorem}[Derived additive KZ connection; \ClaimStatusProvedHere]
\label{thm:derived-additive-kz}
Let $(V,\Omega)$ be a bar-horizontal strictification. On the ordered
configuration space
\[
\Conf_n^{\mathrm{ord}}(\mathbb A^1)=\{z_1<\cdots<z_n\}
\]
consider the operator
\[
\mathbb D_n
=
 d_V+d_{\mathrm{dR}}
 - \sum_{1\le i<j\le n}\Omega_{ij}(z_j-z_i)\,d(z_j-z_i).
\]
Then $\mathbb D_n^2=0$.
\end{theorem}

\begin{proof}
Write
\[
A_n:=\sum_{i<j}\Omega_{ij}(z_j-z_i)\,d(z_j-z_i).
\]
Since $d_V^2=d_{\mathrm{dR}}^2=0$ and $[d_V,\Omega_{ij}(u)]=0$, one has
\[
(d_V+d_{\mathrm{dR}})^2=0,
\qquad
[d_V+d_{\mathrm{dR}},A_n]=0,
\]
because $d_{\mathrm{dR}}(f(u)du)=f'(u)du\wedge du=0$.
Hence
\[
\mathbb D_n^2=A_n\wedge A_n.
\]
Disjoint-pair terms vanish by hypothesis. Thus only triples
$i<j<k$ contribute. Put
\[
u=z_j-z_i,
\qquad
v=z_k-z_j,
\qquad
z_k-z_i=u+v.
\]
Then
\[
d(z_k-z_i)=du+dv,
\]
and every wedge product appearing from the triple $i<j<k$ reduces to a
multiple of $du\wedge dv$. The coefficient is precisely
\[
[\Omega_{ij}(u),\Omega_{ik}(u+v)]
+[\Omega_{ij}(u),\Omega_{jk}(v)]
+[\Omega_{ik}(u+v),\Omega_{jk}(v)],
\]
which vanishes by the defining spectral Kohno relation. Therefore
$\mathbb D_n^2=0$.
\end{proof}

\begin{corollary}[Derived Drinfeld--Kohno shadow; \ClaimStatusProvedHere]
If the filtration is pronilpotent, the parallel transport of
$\mathbb D_n$ is well-defined and yields chain-level pure braid group
representations on $V^{\widehat\otimes n}$, hence braid-group actions
on $H^\bullet(V^{\widehat\otimes n})$.
\end{corollary}

\begin{theorem}[Boundary residue theorem; \ClaimStatusProvedHere]
\label{thm:boundary-residue}
Assume that each $\Omega_{ij}(u)$ has at most a simple pole at $u=0$,
so that
\[
\Omega_{ij}(u)=\frac{r_{ij}}{u}+\Omega_{ij}^{\mathrm{reg}}(u).
\]
Then the connection form of Theorem~\ref{thm:derived-additive-kz}
extends logarithmically to the Fulton--MacPherson compactification
$\overline{\Conf}_n^{\mathrm{ord}}(\mathbb A^1)$.
For the boundary divisor $D_S$ corresponding to collision of a subset
$S\subset\{1,\dots,n\}$, its residue is
\[
\operatorname{Res}_{D_S}(A_n)=\sum_{i<j\in S} r_{ij}.
\]
\end{theorem}

\begin{proof}
Near the divisor $D_S$, choose local coordinates
\[
z_i=z_*+t\xi_i\qquad (i\in S),
\]
with $t=0$ cutting out the boundary and the $\xi_i$ pairwise distinct.
Then for $i,j\in S$,
\[
z_j-z_i=t(\xi_j-\xi_i),
\qquad
 d(z_j-z_i)=(\xi_j-\xi_i)dt+t\,d(\xi_j-\xi_i).
\]
Hence
\begin{align*}
\Omega_{ij}(z_j-z_i)\,d(z_j-z_i)
&=
\left(\frac{r_{ij}}{t(\xi_j-\xi_i)}+\text{regular}\right)
\big((\xi_j-\xi_i)dt+t\,d(\xi_j-\xi_i)\big)\\
&=r_{ij}\frac{dt}{t}+r_{ij}d\log(\xi_j-\xi_i)+\text{regular}.
\end{align*}
Pairs not entirely contained in $S$ remain regular in $t$.
Summing over $i<j\in S$ gives the residue formula.
\end{proof}

\begin{theorem}[Transfer of flat spectral connections; \ClaimStatusProvedHere]
\label{thm:transfer-flat-spectral}
Let
\[
(H,d_H)
\xrightarrow{i}
(V,d_V)
\xrightarrow{p}
(H,d_H)
\]
be a filtered strong deformation retract with homotopy $h$:
\[
pi=\operatorname{id}_H,
\qquad
ip=\operatorname{id}_V+d_Vh+hd_V.
\]
Assume that $V\otimes \Omega^\bullet(\Conf_n^{\mathrm{ord}})$ carries a
flat total differential
\[
\mathbb D_V=(d_V+d_{\mathrm{dR}})-A_V,
\]
where $A_V$ lowers the filtration. Then the homological perturbation
series converges and produces a flat total differential
\[
\mathbb D_H=(d_H+d_{\mathrm{dR}})-A_H,
\qquad
A_H=pA_V(1-hA_V)^{-1}i=
\sum_{m\ge 0} pA_V(hA_V)^mi.
\]
Thus every flat bar-level spectral connection descends functorially to
cohomology or to a minimal model.
\end{theorem}

\begin{proof}
Tensor the strong deformation retract with the de Rham complex of the
base. Since $d_{\mathrm{dR}}$ commutes with $i,p,h$, the same data form a
strong deformation retract for the base differential
$d_V+d_{\mathrm{dR}}$. Set $\delta=-A_V$. Because $\delta$ lowers the
filtration, the standard homological perturbation series converges.
The perturbation lemma gives the transferred differential
\[
(d_H+d_{\mathrm{dR}})+p\delta(1-h\delta)^{-1}i,
\]
which squares to zero because
$(d_V+d_{\mathrm{dR}}+\delta)^2=\mathbb D_V^2=0$.
Replacing $\delta$ by $-A_V$ yields the formula for $A_H$.
\end{proof}

\section{Quasi-factorization and spectral obstruction theory}

\begin{definition}[Regularized holonomy data]
Assume that the flat connections of
Theorem~\ref{thm:derived-additive-kz} are pronilpotent at the boundary.
For $n=2$, let
\[
R(u)\in 1+F^1A^{\widehat\otimes 2}
\]
be the regularized holonomy from the positive tangential basepoint at
collision to the interior point with difference $u=z_2-z_1$.
For $n=3$, let $\Phi(w,z)\in 1+F^2A^{\widehat\otimes 3}$ be the
regularized parallel transport between the two tangential basepoints
corresponding to the bracketings $(12)3$ and $1(23)$ in
$\overline{\Conf}_3^{\mathrm{ord}}(\mathbb A^1)$.
We call $(R,\Phi)$ the \emph{spectral quasi-factorization datum}.
\end{definition}

\begin{theorem}[Quasi-factorization theorem; \ClaimStatusProvedHere]
\label{thm:quasi-factorization}
Every residue-bounded complete dg-shifted Yangian equipped with a
bar-horizontal strictification canonically determines a spectral
quasi-factorization quantum group.
More precisely, the regularized holonomy data $(R,\Phi)$ satisfy:
\begin{enumerate}
\item quasi-hexagon relations on $\overline{\Conf}_3^{\mathrm{ord}}$, and
\item the pentagon relation on the codimension-one Stasheff cell complex
 of $\overline{\Conf}_4^{\mathrm{ord}}$.
\end{enumerate}
\end{theorem}

\begin{proof}
The codimension-one boundary strata of
$\overline{\Conf}_3^{\mathrm{ord}}$ corresponding to the bracketings
$(12)3$ and $1(23)$ admit collar charts in which the regularized
holonomy factors as the ordered product of the two-body holonomies.
The comparison between these two collar factorizations is exactly the
comparison holonomy $\Phi(w,z)$, which yields the quasi-hexagon.

For $n=4$, the codimension-one boundary strata form the Stasheff
pentagon. Flatness of the four-point connection implies that the total
regularized transport around the boundary of that pentagon is the
identity. Writing that transport as the product of the five comparison
holonomies gives the pentagon relation.
\end{proof}

\begin{definition}[Spectral cochain complex and spectral Drinfeld class]
Let
\[
C^2_{\mathrm{spec}}(A)=\{g(u)\in A^{\widehat\otimes 2}\},
\qquad
C^3_{\mathrm{spec}}(A)=\{\phi(w,z)\in A^{\widehat\otimes 3}\}.
\]
Define the linearized spectral differential by
\[
(\delta g)(w,z)
=
1\otimes g(z)
+ (\operatorname{id}\otimes \Delta_z)g(w)
- (\Delta_w\otimes \operatorname{id})g(z)
- g(w)\otimes 1.
\]
The degree-three cohomology of the resulting complex is denoted by
$H^3_{\mathrm{spec}}(A)$.
If
\[
\Phi(w,z)=1+\phi_n(w,z)\pmod{F^{n+1}}
\]
with $\phi_n$ the first nontrivial term, then its class
\[
\mathfrak D_{\mathrm{spec}}(A,\Omega):=[\phi_n]
\in H^3_{\mathrm{spec}}(\operatorname{gr}^n A)
\]
is called the \emph{spectral Drinfeld class}.
\end{definition}

\begin{remark}[Well-definedness of the spectral Drinfeld class]
The spectral Drinfeld class is well-defined because the Hochschild-type
obstruction cocycle it measures is independent of the choice of
filtration splitting: different splittings produce cohomologous cocycles,
so the class in $H^3_{\mathrm{spec}}(\operatorname{gr}^n A)$ is intrinsic.
\end{remark}

\begin{theorem}[Strictification by spectral cohomology; \ClaimStatusProvedHere]
\label{thm:strictification-spectral-cohomology}
Suppose
\[
\Phi(w,z)=1+\phi_n(w,z)\pmod{F^{n+1}}
\]
with $n\ge 2$. Then:
\begin{enumerate}
\item $\phi_n$ is a $3$-cocycle in the spectral complex;
\item under a spectral twist $G_n(u)=1+g_n(u)\pmod{F^{n+1}}$ one has
 \[
 \phi_n\longmapsto \phi_n+\delta g_n;
 \]
\item if $[\phi_n]=0$ in
 $H^3_{\mathrm{spec}}(\operatorname{gr}^nA)$, then a twist exists which
 kills $\phi_n$ modulo $F^{n+1}$;
\item if $H^3_{\mathrm{spec}}(\operatorname{gr}^nA)=0$ for all $n\ge 2$,
 then there is a filtered twist $G$ for which $\Phi^G=1$, hence the
 quasi-factorization datum strictifies to an honest spectral
 factorization $R$-matrix.
\end{enumerate}
\end{theorem}

\begin{proof}
Expand the pentagon relation in the first nontrivial filtration degree.
All nonlinear terms are of higher filtration, so the degree-$n$ term
satisfies the linearized pentagon, i.e. it is a cocycle. Expanding the
standard twist formula for the quasi-associator to degree $n$ yields the
addition by the coboundary $\delta g_n$. Killing the first nontrivial
term and iterating recursively proves the last two statements.
\end{proof}

\begin{conclusion}[First conceptual reduction]
The dg-shifted Yangian/factorization bridge factors through two steps.
First the canonical passage from bar-horizontal strictification data
for a dg-shifted Yangian to a spectral quasi-factorization package,
\[
\text{dg-shifted Yangian + bar-horizontal strictification}
\longrightarrow
\text{spectral quasi-factorization quantum group};
\]
then the strictification problem, controlled by a concrete filtered
cohomology theory.
\end{conclusion}

\section{Low-order algebra: quadratic and cubic control}

\begin{proposition}[Quadratic hexagon obstruction; \ClaimStatusProvedHere]
Let $\widetilde r(z)$ be the suspended degree-zero kernel. Define the
naive multiplicative two-jet
\[
R^{(2)}(z)=1+\widetilde r(z)+\frac12\widetilde r(z)^{\star 2}.
\]
Assume the infinitesimal coassociativity identities
\begin{align*}
(\operatorname{id}\otimes \Delta_z)\widetilde r(z+w)
&=\widetilde r_{12}(w)+\widetilde r_{13}(z+w),\\
(\Delta_w\otimes \operatorname{id})\widetilde r(z)
&=\widetilde r_{13}(z+w)+\widetilde r_{23}(z).
\end{align*}
Then modulo $F^3$,
\begin{align*}
(\operatorname{id}\otimes \Delta_z)R^{(2)}(z+w)-R^{(2)}_{12}(w)R^{(2)}_{13}(z+w)
&\equiv \frac12[\widetilde r_{13}(z+w),\widetilde r_{12}(w)],\\
(\Delta_w\otimes \operatorname{id})R^{(2)}(z)-R^{(2)}_{13}(z+w)R^{(2)}_{23}(z)
&\equiv \frac12[\widetilde r_{23}(z),\widetilde r_{13}(z+w)].
\end{align*}
\end{proposition}

\begin{proof}
Expand both sides to second order and subtract. The square terms cancel
except for the order-sensitive cross terms, which produce the displayed
commutators.
\end{proof}

\begin{proposition}[Three-particle obstruction identity; \ClaimStatusProvedHere]
Let $B^{\le 3}(Y)$ be the reduced bar coalgebra truncated at tensor
length $3$, and let $I_{\widetilde r}$ be the adjacent-gap insertion
coderivation determined by $\widetilde r$. Put
\[
D=d_B+I_{\widetilde r}.
\]
Assume the one-body and two-body Maurer--Cartan identities already
hold. Then on tensor length $3$,
\[
D^2=I_{\mathfrak Y_{\widetilde r}},
\]
where
\[
\mathfrak Y_{\widetilde r}(z,w)=
MC(\widetilde r_{12}(w)+\widetilde r_{13}(z+w)+\widetilde r_{23}(z))
-MC(\widetilde r_{12}(w))
-MC(\widetilde r_{13}(z+w))
-MC(\widetilde r_{23}(z)).
\]
Thus the tensor-length-$3$ obstruction is exactly the
$A_\infty$ Yang--Baxter defect.
\end{proposition}

\begin{proof}
Expand
\[
D^2=[d_B,I_{\widetilde r}] + I_{\widetilde r}^2.
\]
On length $3$, disjoint insertions cancel by coderivation signs.
The surviving terms are exactly the overlapping insertions in the three
ordered gaps. These assemble into the bar realization of the displayed
Maurer--Cartan defect.
\end{proof}

\begin{theorem}[Abelian strictification theorem; \ClaimStatusProvedHere]
\label{thm:abelian-strictification}
Assume that all coefficients of $\widetilde r(z)$ supercommute after
suspension, that $m_k=0$ for $k\ge 3$, and that the underlying
coproduct is cocommutative. Then
\[
R(z)=\exp_\star(\widetilde r(z))
\]
converges and defines an honest spectral $R$-matrix on the completed
bialgebra underlying $Y$.
\end{theorem}

\begin{proof}
The infinitesimal coassociativity identities exponentiate because all
coefficients commute. The hexagon identities for the exponentiated
coproduct follow from the BCH formula: since all coefficients
commute, BCH reduces to ordinary addition, and the hexagon relation
exponentiates termwise. Cocommutativity implies almost
cocommutativity automatically in the commuting completion.
\end{proof}

\begin{theorem}[Residue-bounded completion theorem; \ClaimStatusProvedHere]
\label{thm:residue-bounded-completion}
Let $Y=\varprojlim_N Y_N$ be residue-bounded and complete, where
$Y_N=Y/F^{N+1}Y$. Assume:
\begin{enumerate}
\item each $Y_N$ is quadratic-Koszul;
\item the induced kernels satisfy the three-particle identity above;
\item the inverse system $H^m(B(Y_N))$ is Mittag--Leffler for every $m$.
\end{enumerate}
Then the completed coderivation
\[
D=d_B+I_{\widetilde r}
\]
converges continuously on $\widehat B(Y)$, satisfies $D^2=0$, and the
completed cobar comparison
\[
\Omega\widehat B(Y)\longrightarrow Y
\]
is a quasi-isomorphism.
\end{theorem}

\begin{proof}
Residue-boundedness implies that on each fixed filtration quotient only
finitely many residues contribute, hence the coderivation converges.
The three-particle obstruction identity implies $D^2=0$ quotientwise;
separatedness of the filtration ($\bigcap_n F^n = 0$ by the
pole-order bound) passes the identity to the inverse limit.
Quadratic-Koszulness gives quotientwise bar--cobar
quasi-isomorphisms (by \cite[Proposition~2.2.1]{LV12}), and the
Mittag--Leffler hypothesis (which holds because the transition maps
$H^m(B(Y_{N+1}))\to H^m(B(Y_N))$ are surjective) passes
quasi-isomorphism to the inverse limit.
\end{proof}

\section{Exact Cartan transport and the first genuine shifted examples}

\begin{construction}[Cartan shift cocycle]
Let $H$ be a completed commutative primitive Cartan bialgebra with
primitive generators $h_1,\dots,h_\ell$. Suppose the shifted Cartan
currents are
\[
H_i(u)=
\frac{\prod_j (u-a_j^{(i)})}{\prod_k (u-b_k^{(i)})}.
\]
Set
\[
\phi_i(u)=\log H_i(u).
\]
Define the formal bicharacter
\[
F_\mu(u)=
\exp\!\left(\sum_{i=1}^{\ell} \phi_i(u)\,h_i\otimes h_i\right).
\]
\end{construction}

\begin{proposition}[Cartan shift cocycle identities; \ClaimStatusProvedHere]
\label{prop:cartan-shift-cocycle}
The formal bicharacter $F_\mu(u)$ satisfies
\begin{align*}
(\operatorname{id}\otimes \Delta)F_\mu(u)
&=F_\mu^{12}(u)F_\mu^{13}(u),\\
(\Delta\otimes \operatorname{id})F_\mu(u)
&=F_\mu^{13}(u)F_\mu^{23}(u),
\end{align*}
and is additive in the shift parameter whenever the logarithms are.
\end{proposition}

\begin{proof}
Since $\Delta h_i=h_i\otimes 1+1\otimes h_i$ and the $h_i$ commute,
the exponent splits exactly under either iterated coproduct.
\end{proof}

\begin{theorem}[Cartan gauge-twist theorem; \ClaimStatusProvedHere]
\label{thm:cartan-gauge-twist}
Let
\[
\mathbb D_n^0=(d+d_{\mathrm{dR}})-A_n^0
\]
be an unshifted derived KZ differential. Suppose that on the chosen
sector the operators $h_i^{(a)}h_i^{(b)}$ commute with $A_n^0$.
Set
\[
S_{\mu,n}=
\sum_{1\le a<b\le n}\sum_{i=1}^{\ell}
\phi_i(z_b-z_a)\,h_i^{(a)}h_i^{(b)},
\qquad
G_{\mu,n}=e^{-S_{\mu,n}}.
\]
Then
\[
\mathbb D_n^\mu=G_{\mu,n}\,\mathbb D_n^0\,G_{\mu,n}^{-1}
=(d+d_{\mathrm{dR}})-A_n^\mu
\]
has connection form
\[
A_n^\mu
=
A_n^0
+
\sum_{1\le a<b\le n}\sum_{i=1}^{\ell}
\phi_i'(z_b-z_a)\,h_i^{(a)}h_i^{(b)}\,d(z_b-z_a).
\]
\end{theorem}

\begin{proof}
The $h_i$ commute in the dg-shifted setting because the Cartan
subalgebra $\mathfrak{h}\subset\mathfrak{g}$ is abelian, and the
dg-shifted coproduct preserves commutativity. Because the exponents
commute on the sector in question,
\[
dG_{\mu,n}\,G_{\mu,n}^{-1}=-dS_{\mu,n}.
\]
The stated commutation hypothesis implies
$G_{\mu,n}A_n^0G_{\mu,n}^{-1}=A_n^0$, so the gauge transform contributes
exactly $dS_{\mu,n}$.
\end{proof}

\begin{theorem}[Cartan-diagonal shift defect is exact; \ClaimStatusProvedHere]
\label{thm:cartan-diagonal-defect-exact}
Fix a Cartan-diagonal line $L_\lambda$ and write
\[
\phi_i(u)=\log H_i(u).
\]
The shifted two-body scalar factor is
\[
R_{\mu,\lambda}^{\mathrm{sh}}(u)=
\exp\!\Big(\sum_i \lambda_i^2\phi_i(u)\Big).
\]
The three-body quasi-hexagon defect is therefore
\[
\Phi_{\mu,\lambda}^{\mathrm{sh}}(w,z)=
\frac{R_{\mu,\lambda}^{\mathrm{sh}}(w+z)}{R_{\mu,\lambda}^{\mathrm{sh}}(w)
R_{\mu,\lambda}^{\mathrm{sh}}(z)},
\]
and its logarithm is
\[
\mathfrak D_{\mathrm{spec},\mu}^{\lambda}(w,z)=
\sum_i \lambda_i^2\big(\phi_i(w+z)-\phi_i(w)-\phi_i(z)\big).
\]
Hence the Cartan-diagonal contribution to the spectral Drinfeld class is
exact.
\end{theorem}

\begin{proof}
Immediate from the definition of the shifted scalar transport.
\end{proof}

\begin{remark}[First genuine shifted examples]
The framed $A_1$ shifted Yangian and the staircase shifted Yangians of
finite type~$A$ furnish the first genuine shifted examples. In all such
Cartan-diagonal sectors, the shift produces a visible three-body
logarithmic defect but no cohomological obstruction class.
\end{remark}

\section{Filtration $2$: adjacent-root triangles}

The first genuinely mixed-root sector arises from adjacent simple roots.
Let $i\sim j$ be adjacent simple roots in a rank-$2$ Levi of type $A_2$.
Set
\[
E_{ij}=[E_i,E_j],
\qquad
F_{ij}=[F_j,F_i],
\]
so that $[F_i,F_j]=-F_{ij}$.
For $a\neq b$, write
\[
X_i^{ab}(u)=\rho_i(u)E_i^{(a)}F_i^{(b)},
\qquad
X_{ij}^{ab}(u)=\rho_{ij}(u)E_{ij}^{(a)}F_{ij}^{(b)}.
\]

\begin{definition}[Triangle sectors]
On three tensor factors define the left and right adjacent-root triangle
vectors
\[
T_{ij}^{L}=E_{ij}^{(1)}F_i^{(2)}F_j^{(3)},
\qquad
T_{ij}^{R}=E_i^{(1)}E_j^{(2)}F_{ij}^{(3)}.
\]
The corresponding opposite orientations are
\[
T_{ji}^{L}=E_{ij}^{(1)}F_j^{(2)}F_i^{(3)},
\qquad
T_{ji}^{R}=E_j^{(1)}E_i^{(2)}F_{ij}^{(3)}.
\]
\end{definition}

\begin{theorem}[Triangle localization formula; \ClaimStatusProvedHere]
\label{thm:triangle-localization}
Let
\[
S_{12}(w)=\rho_i(w)E_i^{(1)}F_i^{(2)}+\rho_j(w)E_j^{(1)}F_j^{(2)}+\cdots,
\]
\[
S_{13}(w+z)=\rho_i(w+z)E_i^{(1)}F_i^{(3)}
+\rho_j(w+z)E_j^{(1)}F_j^{(3)}+\cdots.
\]
Then
\[
\Pi_{\langle T_{ij}^{L},T_{ji}^{L}\rangle}\frac12[S_{12}(w),S_{13}(w+z)]
=
\frac12\rho_i(w)\rho_j(w+z)T_{ij}^{L}
-
\frac12\rho_j(w)\rho_i(w+z)T_{ji}^{L}.
\]
Similarly,
\[
\Pi_{\langle T_{ij}^{R},T_{ji}^{R}\rangle}\frac12[S_{13}(w+z),S_{23}(z)]
=
-\frac12\rho_i(w+z)\rho_j(z)T_{ij}^{R}
+
\frac12\rho_j(w+z)\rho_i(z)T_{ji}^{R}.
\]
\end{theorem}

\begin{proof}
The only contributing commutators are
\[
[E_i^{(1)}F_i^{(2)},E_j^{(1)}F_j^{(3)}]
=E_{ij}^{(1)}F_i^{(2)}F_j^{(3)},
\]
and
\[
[E_j^{(1)}F_j^{(2)},E_i^{(1)}F_i^{(3)}]
=-E_{ij}^{(1)}F_j^{(2)}F_i^{(3)},
\]
with the analogous computation on the right using
$[F_i,F_j]=-F_{ij}$.
\end{proof}

\begin{definition}[Triangle coproduct coefficients]
Let $\Delta_z$ be the twisted coproduct. Define the triangle transport
coefficients by
\[
\Pi_{\langle T_{ij}^{L},T_{ji}^{L}\rangle}
\big((\operatorname{id}\otimes\Delta_z)X_{ij}^{12}(w)\big)
=
\kappa_{ij}^{L}(w,z)T_{ij}^{L}-\kappa_{ji}^{L}(w,z)T_{ji}^{L},
\]
and similarly on the right,
\[
\Pi_{\langle T_{ij}^{R},T_{ji}^{R}\rangle}
\big((\Delta_w\otimes\operatorname{id})X_{ij}^{23}(z)\big)
=
-\kappa_{ij}^{R}(w,z)T_{ij}^{R}+\kappa_{ji}^{R}(w,z)T_{ji}^{R}.
\]
\end{definition}

\begin{theorem}[Adjacent-root rigidity; \ClaimStatusProvedHere]
\label{thm:adjacent-root-rigidity}
The filtration-$2$ hexagon defect vanishes in the oriented triangle
sectors if and only if
\begin{align*}
\rho_{ij}(w)\kappa_{ij}^{L}(w,z)&=\frac12\rho_i(w)\rho_j(w+z),\\
\rho_{ij}(w)\kappa_{ji}^{L}(w,z)&=\frac12\rho_j(w)\rho_i(w+z),\\
\rho_{ij}(z)\kappa_{ij}^{R}(w,z)&=\frac12\rho_i(w+z)\rho_j(z),\\
\rho_{ij}(z)\kappa_{ji}^{R}(w,z)&=\frac12\rho_j(w+z)\rho_i(z).
\end{align*}
\end{theorem}

\begin{proof}
Project the logarithm of the product side to the triangle sectors.
The only contributions are the commutator terms computed in
Theorem~\ref{thm:triangle-localization}. On the coproduct side,
$X_{ij}$ is the unique root generator whose twisted coproduct can land in
those sectors. Equality of coefficients gives the claim.
\end{proof}

\begin{corollary}[Normalized rational triangle coefficient; \ClaimStatusProvedHere]
Assume the rational ansatz
\[
\rho_i(u)=\frac{c_i}{u},
\qquad
\rho_j(u)=\frac{c_j}{u},
\qquad
\rho_{ij}(u)=\frac{c_{ij}}{u},
\]
and
\[
\kappa_{ij}^{L}(w,z)=\frac{\lambda_{ij}^{L}}{w+z},
\qquad
\kappa_{ij}^{R}(w,z)=\frac{\lambda_{ij}^{R}}{w+z},
\]
with the same form for the opposite orientations. Then
\[
\lambda_{ij}^{L}=
\lambda_{ji}^{L}=
\lambda_{ij}^{R}=
\lambda_{ji}^{R}=
\frac{c_ic_j}{2c_{ij}}.
\]
In the normalized case $c_i=c_j=c_{ij}=1$, the universal coefficient is
$\frac12$.
\end{corollary}

\begin{proof}
Substitute the rational ansatz into the equations of
Theorem~\ref{thm:adjacent-root-rigidity}.
\end{proof}

\begin{remark}[Transport, not coefficientwise invariance]
\label{rem:transport-not-invariance}
At mixed-root level the shift changes the local mixed-root coefficients
by explicit ratios of shift characters at \emph{different spectral
arguments}. The invariant is the obstruction class after transport by
the shift cocycle, not the raw coefficient.
\end{remark}

\begin{proposition}[Triangle shift transport law; \ClaimStatusProvedHere]
\label{prop:triangle-shift-transport}
Suppose the shifted root coefficients satisfy
\[
\rho_\alpha^{\mu}(u)=\chi_\alpha^{\mu}(u)\,\rho_\alpha^{0}(u)
\]
for scalar shift characters $\chi_\alpha^{\mu}(u)$.
Assume the unshifted triangle coefficients solve the rigidity equations.
Then the shifted coefficients are uniquely forced to be
\begin{align*}
\kappa_{ij}^{L,\mu}(w,z)
&=
\frac{\chi_i^{\mu}(w)\chi_j^{\mu}(w+z)}{\chi_{ij}^{\mu}(w)}
\,\kappa_{ij}^{L,0}(w,z),\\
\kappa_{ij}^{R,\mu}(w,z)
&=
\frac{\chi_i^{\mu}(w+z)\chi_j^{\mu}(z)}{\chi_{ij}^{\mu}(z)}
\,\kappa_{ij}^{R,0}(w,z),
\end{align*}
with the analogous formulas for the opposite orientations.
\end{proposition}

\begin{proof}
Insert the shifted coefficients into the identities of
Theorem~\ref{thm:adjacent-root-rigidity} and solve for
$\kappa^{L,\mu}$ and $\kappa^{R,\mu}$.
\end{proof}

\begin{interpretation}
The first mixed-root sector detects shift sensitivity at the level of
local transport functions. This sensitivity is forced by the shifted
one-body characters; no new independent obstruction class arises.
\end{interpretation}

\section{The exact free two-body model}

The free example of \cite{DNP25} admits a particularly transparent
computation. Write
\[
X(x)=\sum_{n\ge 0} X_n x^n,
\qquad
\Psi(x)=\sum_{n\ge 0} \psi_n x^n.
\]
The dg kernel in the free model is
\[
r(u)=\sum_{n,m\ge 0}(-1)^n\binom{n+m}{m}
\frac{\psi_n\otimes X_m-X_n\otimes \psi_m}{u^{n+m+1}}.
\]

\begin{lemma}[Closed generating kernel; \ClaimStatusProvedHere]
\label{lem:free-closed-generating-kernel}
Coefficientwise in $x,y$,
\[
\sum_{n,m\ge 0}(-1)^n\binom{n+m}{m}\frac{x^n y^m}{u^{n+m+1}}
=
\frac{1}{u+x-y}.
\]
Hence the free suspended kernel is the coefficientwise expansion of
\[
\mathcal R_{\mathrm{free}}(u;x,y)=
\frac{\Psi(x)\otimes X(y)-X(x)\otimes \Psi(y)}{u+x-y}.
\]
\end{lemma}

\begin{proof}
Summing over $N=n+m$ gives
\begin{align*}
\sum_{n,m\ge 0}(-1)^n\binom{n+m}{m}\frac{x^n y^m}{u^{n+m+1}}
&=
\frac1u\sum_{N\ge 0}\frac1{u^N}
\sum_{m=0}^N \binom{N}{m}(-x)^{N-m}y^m\\
&=
\frac1u\sum_{N\ge 0}\left(\frac{y-x}{u}\right)^N
=
\frac{1}{u+x-y}.
\end{align*}
\end{proof}

\begin{theorem}[Exact free two-body transport; \ClaimStatusProvedHere]
\label{thm:exact-free-transport}
Let
\[
\Theta(x,y)=\Psi(x)\otimes X(y)-X(x)\otimes \Psi(y).
\]
In the free supercommutative model, $\Theta(x,y)^2=0$. Therefore the
connection
\[
\nabla_{\mathrm{free}}=d-
\frac{\Theta(x,y)}{u+x-y}\,du
\]
integrates exactly with transport
\[
T(u,u_0)=1+\Theta(x,y)\log\!\frac{u+x-y}{u_0+x-y}.
\]
The local monodromy around the formal divisor $u+x-y=0$ is
\[
M_{\mathrm{loop}}=1+2\pi i\,\Theta(x,y).
\]
\end{theorem}

\begin{proof}
Because $\Theta^2=0$,
\[
\exp\left(\Theta\log\!\frac{u+x-y}{u_0+x-y}\right)
=1+\Theta\log\!\frac{u+x-y}{u_0+x-y}.
\]
Differentiating gives the transport equation. The monodromy follows by
adding $2\pi i$ to the logarithm around a loop.
\end{proof}

\section{Filtration $3$: root quadrilaterals}

The first class-$3$ computation.
Let $i\sim j\sim k$ be a chain of simple roots of type $A_3$ with
$i$ not adjacent to $k$. Define the composite roots by
\[
E_{ij}=[E_i,E_j],
\qquad
E_{jk}=[E_j,E_k],
\qquad
E_{ijk}=[E_{ij},E_k]=[E_i,E_{jk}],
\]
and for the opposite side set
\[
F_{ij}=[F_j,F_i],
\qquad
F_{jk}=[F_k,F_j],
\qquad
F_{ijk}=[F_k,F_{ij}]=[F_{jk},F_i].
\]
Then
\[
[F_i,F_j]=-F_{ij},
\qquad
[F_j,F_k]=-F_{jk},
\qquad
[F_i,F_{jk}]=-[F_k,F_{ij}]=-F_{ijk}.
\]

\begin{definition}[Root-quadrilateral sectors]
On four tensor factors define the oriented left and right quadrilateral
vectors
\[
Q_{ijk}^{L}=E_{ijk}^{(1)}F_i^{(2)}F_j^{(3)}F_k^{(4)},
\qquad
Q_{ijk}^{R}=E_i^{(1)}E_j^{(2)}E_k^{(3)}F_{ijk}^{(4)}.
\]
The opposite orientation is obtained by replacing $(i,j,k)$ by
$(k,j,i)$ and the sign conventions above.
\end{definition}

\begin{lemma}[Class-$3$ ordered BCH coefficient; \ClaimStatusProvedHere]
\label{lem:class3-ordered-bch}
Let $X,Y,Z$ lie in a class-$3$ nilpotent Lie algebra, so all
commutators of depth $4$ vanish. The multilinear $XYZ$-part of
$\log(e^Xe^Ye^Z)$ is
\[
\frac14[[X,Y],Z]+\frac1{12}[X,[Y,Z]]+\frac1{12}[Y,[X,Z]].
\]
If, in addition, $[X,Z]=0$, then this simplifies to
\[
\operatorname{mult}_{XYZ}\log(e^Xe^Ye^Z)
=\frac13[X,[Y,Z]].
\]
\end{lemma}

\begin{proof}
Write
\[
U=\log(e^Xe^Y)=X+Y+\frac12[X,Y]+\frac1{12}[X,[X,Y]]
+\frac1{12}[Y,[Y,X]].
\]
Now compute
\[
\log(e^Xe^Ye^Z)=\log(e^Ue^Z)
=U+Z+\frac12[U,Z]+\frac1{12}[U,[U,Z]]+\frac1{12}[Z,[Z,U]].
\]
Retaining only the multilinear $XYZ$-part gives the first displayed
formula. If $[X,Z]=0$, then
\[
[[X,Y],Z]=[X,[Y,Z]]
\]
by Jacobi, and the second formula follows.
\end{proof}

\begin{theorem}[Quadrilateral localization formula; \ClaimStatusProvedHere]
\label{thm:quadrilateral-localization}
Let
\[
X_i^{12}(u)=\rho_i(u)E_i^{(1)}F_i^{(2)},
\qquad
X_j^{13}(u+v)=\rho_j(u+v)E_j^{(1)}F_j^{(3)},
\qquad
X_k^{14}(u+v+t)=\rho_k(u+v+t)E_k^{(1)}F_k^{(4)}.
\]
Then the projection of
\[
\log\Big(e^{X_i^{12}(u)}e^{X_j^{13}(u+v)}e^{X_k^{14}(u+v+t)}\Big)
\]
to the left quadrilateral sector is
\[
\frac13\rho_i(u)\rho_j(u+v)\rho_k(u+v+t)\,Q_{ijk}^{L}.
\]
Similarly, the projection of
\[
\log\Big(e^{\rho_i(u+v+t)E_i^{(1)}F_i^{(4)}}
 e^{\rho_j(v+t)E_j^{(2)}F_j^{(4)}}
 e^{\rho_k(t)E_k^{(3)}F_k^{(4)}}\Big)
\]
to the right quadrilateral sector is
\[
-\frac13\rho_i(u+v+t)\rho_j(v+t)\rho_k(t)\,Q_{ijk}^{R}.
\]
\end{theorem}

\begin{proof}
For the left sector, set
\[
X=X_i^{12}(u),\qquad Y=X_j^{13}(u+v),\qquad Z=X_k^{14}(u+v+t).
\]
Since $i$ and $k$ are nonadjacent, $[X,Z]=0$.
By Lemma~\ref{lem:class3-ordered-bch}, the multilinear part is
$\frac13[X,[Y,Z]]$. Now
\[
[Y,Z]=\rho_j(u+v)\rho_k(u+v+t)E_{jk}^{(1)}F_j^{(3)}F_k^{(4)},
\]
and then
\[
[X,[Y,Z]]=
\rho_i(u)\rho_j(u+v)\rho_k(u+v+t)
E_{ijk}^{(1)}F_i^{(2)}F_j^{(3)}F_k^{(4)}.
\]
This proves the first formula.

For the right sector, set
\[
X=\rho_i(u+v+t)E_i^{(1)}F_i^{(4)},
\quad
Y=\rho_j(v+t)E_j^{(2)}F_j^{(4)},
\quad
Z=\rho_k(t)E_k^{(3)}F_k^{(4)}.
\]
Again $[X,Z]=0$, so the multilinear part is
$\frac13[X,[Y,Z]]$. Here
\[
[Y,Z]=-
\rho_j(v+t)\rho_k(t)E_j^{(2)}E_k^{(3)}F_{jk}^{(4)},
\]
so
\[
[X,[Y,Z]]=-
\rho_i(u+v+t)\rho_j(v+t)\rho_k(t)
E_i^{(1)}E_j^{(2)}E_k^{(3)}F_{ijk}^{(4)}.
\]
This is the second formula.
\end{proof}

\begin{definition}[Quadrilateral coproduct coefficients]
Let $\Delta^{(2)}_{v,t}$ denote the two-step twisted coproduct
separating one factor into three ordered factors with spectral gaps
$v,t$. Define the left and right quadrilateral transport coefficients
by
\[
\Pi_{\langle Q_{ijk}^{L}\rangle}
\big((\operatorname{id}\otimes\Delta^{(2)}_{v,t})X_{ijk}^{12}(u)\big)
=
\Lambda_{ijk}^{L}(u,v,t)Q_{ijk}^{L},
\]
\[
\Pi_{\langle Q_{ijk}^{R}\rangle}
\big((\Delta^{(2)}_{u,v}\otimes\operatorname{id})X_{ijk}^{34}(t)\big)
=
\Lambda_{ijk}^{R}(u,v,t)Q_{ijk}^{R}.
\]
\end{definition}

\begin{theorem}[Quadrilateral rigidity; \ClaimStatusProvedHere]
\label{thm:quadrilateral-rigidity}
The filtration-$3$ defect vanishes in the oriented quadrilateral sectors
if and only if
\begin{align*}
\rho_{ijk}(u)\Lambda_{ijk}^{L}(u,v,t)
&=
\frac13\rho_i(u)\rho_j(u+v)\rho_k(u+v+t),\\
\rho_{ijk}(t)\Lambda_{ijk}^{R}(u,v,t)
&=
\frac13\rho_i(u+v+t)\rho_j(v+t)\rho_k(t).
\end{align*}
\end{theorem}

\begin{proof}
Project the logarithm of the ordered product of the relevant two-body
factors to the quadrilateral sector. By
Theorem~\ref{thm:quadrilateral-localization}, the product side has the
stated coefficient.
On the coproduct side, a quadrilateral vector with three distinct simple
lowering or raising operators in separate tensor factors can only arise
from the triple-root generator $X_{ijk}$: any simple-root or two-root
composite generator remains supported in at most two root colors and
cannot land in $Q_{ijk}^{L}$ or $Q_{ijk}^{R}$ after projection.
Equality of coefficients yields the claim.
\end{proof}

\begin{corollary}[Normalized rational quadrilateral coefficient; \ClaimStatusProvedHere]
\label{cor:quadrilateral-normalized-rational}
Assume the rational ansatz
\[
\rho_i(u)=\frac{c_i}{u},
\qquad
\rho_j(u)=\frac{c_j}{u},
\qquad
\rho_k(u)=\frac{c_k}{u},
\qquad
\rho_{ijk}(u)=\frac{c_{ijk}}{u}.
\]
Suppose further that the quadrilateral coefficients have translation-
covariant form
\[
\Lambda_{ijk}^{L}(u,v,t)=\frac{\lambda_{ijk}^{L}}{(u+v)(u+v+t)},
\qquad
\Lambda_{ijk}^{R}(u,v,t)=\frac{\lambda_{ijk}^{R}}{(v+t)(u+v+t)}.
\]
Then
\[
\lambda_{ijk}^{L}=\lambda_{ijk}^{R}=\frac{c_ic_jc_k}{3c_{ijk}}.
\]
In the normalized case $c_i=c_j=c_k=c_{ijk}=1$, the universal
quadrilateral coefficient is
\[
\boxed{\frac13.}
\]
\end{corollary}

\begin{proof}
Substitute the rational ansatz into the identities of
Theorem~\ref{thm:quadrilateral-rigidity}.
\end{proof}

\begin{proposition}[Quadrilateral shift transport law; \ClaimStatusProvedHere]
\label{prop:quadrilateral-shift-transport}
Assume again that the shifted one-body coefficients factor as
\[
\rho_\alpha^{\mu}(u)=\chi_\alpha^{\mu}(u)\,\rho_\alpha^0(u).
\]
If the unshifted coefficients solve
Theorem~\ref{thm:quadrilateral-rigidity}, then the shifted quadrilateral
transport coefficients are uniquely forced to be
\begin{align*}
\Lambda_{ijk}^{L,\mu}(u,v,t)
&=
\frac{\chi_i^{\mu}(u)\chi_j^{\mu}(u+v)\chi_k^{\mu}(u+v+t)}
{\chi_{ijk}^{\mu}(u)}
\,\Lambda_{ijk}^{L,0}(u,v,t),\\
\Lambda_{ijk}^{R,\mu}(u,v,t)
&=
\frac{\chi_i^{\mu}(u+v+t)\chi_j^{\mu}(v+t)\chi_k^{\mu}(t)}
{\chi_{ijk}^{\mu}(t)}
\,\Lambda_{ijk}^{R,0}(u,v,t).
\end{align*}
\end{proposition}

\begin{proof}
Insert the shifted coefficients into the two equations of
Theorem~\ref{thm:quadrilateral-rigidity} and solve for
$\Lambda^{L,\mu}$ and $\Lambda^{R,\mu}$.
\end{proof}

\begin{corollary}[No new independent shift obstruction in filtration $3$; \ClaimStatusProvedHere]
\label{cor:no-new-filtration3-shift-obstruction}
A new shift-sensitive \emph{function} appears in filtration $3$ through
the transport laws of
Proposition~\ref{prop:quadrilateral-shift-transport}; this function is
forced by the shifted one-body characters. No new \emph{independent}
obstruction class appears in the oriented root-quadrilateral sector.
\end{corollary}

\begin{proof}
The quadrilateral coefficient is uniquely determined by lower data once
Theorem~\ref{thm:quadrilateral-rigidity} holds. Therefore its shift
variation contributes no new free parameter and no new independent
cohomology class.
\end{proof}

\begin{remark}[First non-path obstruction]
The low-order picture is sharp.
\begin{enumerate}
\item Cartan-diagonal sectors produce exact defects.
\item Commuting sectors strictify directly.
\item Adjacent-root triangles are forced by the universal coefficient
 $\frac12$, transported explicitly by shift characters.
\item Root quadrilaterals are forced by the universal coefficient
 $\frac13$, transported explicitly by shift characters.
\end{enumerate}
The first filtration where an independent obstruction lives is
$\ge 4$, where new Hall words and higher Magnus/BCH terms enter: the
$D_4$ star gives a two-dimensional free multilinear Lie space.
Theorem~\ref{thm:root-space-one-dim} collapses the corresponding finite
root sector to a scalar, and Theorem~\ref{thm:complete-strictification}
identifies the remaining scalar obstruction. The low-order coefficients
$\frac12$ and $\frac13$ are the first instances of the universal
$\frac1n$ in path sectors.
\end{remark}

\section{Conditional global strictification criterion}

\begin{theorem}[Conditional strictification criterion; \ClaimStatusProvedHere]
Assume:
\begin{enumerate}
\item the unshifted dg object strictifies to a factorization quantum
 group with spectral $R$-matrix $R_0$;
\item the shift is implemented by the Cartan cocycle
 of Proposition~\ref{prop:cartan-shift-cocycle};
\item the infinitesimal logarithm of the shifted spectral transport
 agrees with the dg kernel after residue-bounded completion; and
\item the spectral Drinfeld classes vanish.
\end{enumerate}
Then the shifted dg-shifted Yangian admits a completed
$E_1$-chiral/factorization realization whose representation category is
that of an honest spectral factorization quantum group.
\end{theorem}

\begin{proof}
The unshifted strictification gives a strict spectral object.
Transport by the Cartan cocycle gives the shifted quasi-factorization
transport. Vanishing of the spectral Drinfeld class then strictifies the
shifted quasi-object. Residue-bounded completion upgrades finite-stage
compatibility to the full completed statement.
\end{proof}

\section{The strictification theorem}

\subsection{Path-type root sectors}

\begin{definition}[Path-type sector at filtration $n$]
\label{def:path-sector}
Let $\fg$ be a semisimple Lie algebra. A \emph{path-type sector at
filtration~$n$} is determined by a chain of $n$ simple roots
$\alpha_1,\alpha_2,\ldots,\alpha_n$ in the Dynkin diagram such that
$\alpha_i$ is adjacent to $\alpha_{i+1}$ for $1\le i\le n-1$ and
$\alpha_i$ is not adjacent to $\alpha_j$ for $|i-j|\ge 2$. (Such a
chain exists whenever $\fg$ contains a type-$A_n$ subdiagram.)
Define the composite root generators by right-normed iterated brackets:
\[
E_{\alpha_1\cdots\alpha_k}
:=
[E_{\alpha_1},[E_{\alpha_2},\ldots,[E_{\alpha_{k-1}},E_{\alpha_k}]\cdots]]
\qquad (2\le k\le n),
\]
and similarly for the opposite generators $F_{\alpha_1\cdots\alpha_k}$.
The \emph{left $n$-gon vector} and \emph{right $n$-gon vector} on
$n+1$ tensor factors are
\[
P_{\alpha_1\cdots\alpha_n}^{L}
= E_{\alpha_1\cdots\alpha_n}^{(1)}\,
 F_{\alpha_1}^{(2)}\,F_{\alpha_2}^{(3)}\cdots F_{\alpha_n}^{(n+1)},
\qquad
P_{\alpha_1\cdots\alpha_n}^{R}
= E_{\alpha_1}^{(1)}\,E_{\alpha_2}^{(2)}\cdots
 E_{\alpha_n}^{(n)}\,F_{\alpha_1\cdots\alpha_n}^{(n+1)}.
\]
\end{definition}

\begin{remark}
At $n=2$ these recover the triangle vectors $T_{ij}^{L/R}$ of the
adjacent-root computation. At $n=3$ they recover the quadrilateral
vectors $Q_{ijk}^{L/R}$ (relabeling
$(i,j,k)\mapsto(\alpha_1,\alpha_2,\alpha_3)$). The general case is the
$n$-gon sector.
\end{remark}

\subsection{Path one-dimensionality}

\begin{theorem}[Path Lie algebra one-dimensionality; \ClaimStatusProvedHere]
\label{thm:path-one-dim}
Let $L_{\mathrm{path}}(n)$ be the Lie algebra with generators
$X_1,\ldots,X_n$ and relations $[X_i,X_j]=0$ for $|i-j|\ge 2$. Then
the multilinear degree-$n$ component
\[
L_{\mathrm{path}}^n(X_1,\ldots,X_n)
:=
\operatorname{span}\{\text{Lie monomials involving each }X_i\text{ exactly
once}\}
\]
is one-dimensional, generated by the right-normed bracket
\[
[X_1,[X_2,\ldots,[X_{n-1},X_n]\cdots]].
\]
\end{theorem}

\begin{proof}
By induction on $n$.

\medskip\noindent\textbf{Base case} ($n=2$).
The multilinear degree-$2$ space is spanned by $[X_1,X_2]$.

\medskip\noindent\textbf{Inductive step.}
Let $M$ be a multilinear degree-$n$ Lie monomial in $X_1,\ldots,X_n$.
Write $M=[L,R]$ where $L$ involves variables indexed by $S\subset
\{1,\ldots,n\}$ and $R$ involves $T=\{1,\ldots,n\}\setminus S$.

\emph{Step~1: Connectivity.}
For $[L,R]\neq 0$, both $S$ and $T$ must be single connected subpaths
of $1\text{--}2\text{--}\cdots\text{--}n$. The Lie algebra generated by
any disconnected subset of the path generators is a direct sum, and its
multilinear Lie monomials vanish.

Suppose $S\subseteq\{1,\ldots,n\}$ decomposes as $S=S_1\sqcup S_2$
with $\max(S_1)<\min(S_2)$ and $\max(S_1)+1<\min(S_2)$ (a gap in the
path). Then $[X_a,X_c]=0$ for all $a\in S_1$, $c\in S_2$ (since
$|a-c|\ge 2$). The Lie algebra generated by $\{X_s:s\in S\}$ is
therefore the direct sum
$\mathrm{Lie}(\{X_s:s\in S_1\})\oplus\mathrm{Lie}(\{X_s:s\in S_2\})$.
A multilinear Lie monomial involving generators from both summands is
zero (every bracket tree must eventually pair elements from different
summands, producing a commutator that vanishes in a direct sum).
Hence $L=0$ whenever $S$ is disconnected. The same argument
applies to~$T$ and~$R$.

\emph{Step~2: Interval partition.}
From Step~1, the only non-vanishing decompositions are
$S=\{1,\ldots,m\}$, $T=\{m{+}1,\ldots,n\}$ for some $1\le m\le n-1$
(or the reverse, but $[L,R]=-[R,L]$). By the induction hypothesis,
$L$ is proportional to $[X_1,[X_2,\ldots,X_m]\cdots]$ and $R$ is
proportional to $[X_{m+1},[X_{m+2},\ldots,X_n]\cdots]$. So
$M$ is proportional to
\[
\bigl[[X_1,\ldots,X_m],\,[X_{m+1},\ldots,X_n]\bigr]
\]
(right-normed brackets on each side).

\emph{Step~3: Reduction to the standard generator.}
The identity
\[
\bigl[[X_1,\ldots,X_m],\,[X_{m+1},\ldots,X_n]\bigr]
=
[X_1,[X_2,\ldots,[X_{n-1},X_n]\cdots]]
\]
holds in $L_{\mathrm{path}}(n)$ for every $1\le m\le n-1$, by the
Jacobi identity:
\begin{multline*}
\bigl[[X_1,\ldots,X_m],\,[X_{m+1},\ldots,X_n]\bigr]
\;=\;
\bigl[X_1,\,\bigl[[X_2,\ldots,X_m],\,[X_{m+1},\ldots,X_n]\bigr]\bigr]\\
{}-
\bigl[[X_2,\ldots,X_m],\,\bigl[X_1,\,[X_{m+1},\ldots,X_n]\bigr]\bigr].
\end{multline*}
For the second term: $[X_1,X_j]=0$ for all $j\ge 3$, hence
$[X_1,[X_{m+1},\ldots,X_n]]=0$ whenever $m\ge 2$ (the outermost
bracket involves $X_1$ and a monomial on $\{X_{m+1},\ldots,X_n\}$;
by induction on bracketing depth, every leaf of that monomial has
$[X_1,X_j]=0$ for $j\ge m+1\ge 3$).

Thus the second term vanishes for $m\ge 2$, and we have
\[
\bigl[[X_1,\ldots,X_m],\,[X_{m+1},\ldots,X_n]\bigr]
=
\bigl[X_1,\,\bigl[[X_2,\ldots,X_m],\,[X_{m+1},\ldots,X_n]\bigr]\bigr].
\]
By the induction hypothesis applied to the subpath $\{2,\ldots,n\}$
(which has $n-1$ generators satisfying the path relations), the inner
bracket $[[X_2,\ldots,X_m],[X_{m+1},\ldots,X_n]]$ is proportional to
$[X_2,[X_3,\ldots,X_n]\cdots]$. Hence
\[
M \propto [X_1,[X_2,\ldots,[X_{n-1},X_n]\cdots]].
\]
This is a single generator, so $\dim L_{\mathrm{path}}^n=1$.
\end{proof}

\subsection{Universal BCH coefficient via the Dynkin projection}

\begin{theorem}[Dynkin coefficient for the ordered path BCH; \ClaimStatusProvedHere]
\label{thm:dynkin-bch-path}
Let $X_1,\ldots,X_n$ satisfy the path commutation relations
$[X_i,X_j]=0$ for $|i-j|\ge 2$. Then the multilinear
$X_1\cdots X_n$-part of the ordered Baker--Campbell--Hausdorff product is
\begin{equation}\label{eq:bch-path-coeff}
\operatorname{mult}_{X_1\cdots X_n}
\log(e^{X_1}\,e^{X_2}\cdots e^{X_n})
\;=\;
\frac{1}{n}\,[X_1,[X_2,\ldots,[X_{n-1},X_n]\cdots]].
\end{equation}
\end{theorem}

\begin{proof}
By Theorem~\ref{thm:path-one-dim}, the multilinear part is
$c_n\cdot[X_1,[X_2,\ldots,X_n]\cdots]$ for a unique scalar $c_n$. The
scalar $c_n$ is determined by passing to the universal enveloping
algebra and extracting the coefficient of the concatenation monomial
$X_1\cdot X_2\cdots X_n$.

\medskip\noindent\textbf{Step~1: Coefficient in the Lie bracket.}
In the universal enveloping algebra, the right-normed bracket
$[X_1,[X_2,\ldots,[X_{n-1},X_n]\cdots]]$ expands as a signed sum of
permuted monomials. Its leading term (in the order
$X_1X_2\cdots X_n$) has coefficient $+1$.

\medskip\noindent\textbf{Step~2: Coefficient in the BCH logarithm.}
In the free associative algebra, $e^{X_1}\cdot e^{X_2}\cdots e^{X_n}$
has multilinear part exactly $X_1\cdot X_2\cdots X_n$ (each $X_i$
appears in a single exponential factor, so the multilinear extraction
picks $X_i$ from the degree-$1$ term of $e^{X_i}$ and $1$ from all
other factors).

The multilinear part of $\log(P)$ where $P=e^{X_1}\cdots e^{X_n}$ is
computed via $\log(1+A)=\sum_{k\ge 1}(-1)^{k+1}A^k/k$ with $A=P-1$.
The subset-multilinear decomposition of $A$ gives: the
$S$-multilinear part of $A$ (for $S\subseteq\{1,\ldots,n\}$,
$S\neq\emptyset$) is the ordered product $\prod_{i\in S}^{\to}X_i$.
Therefore
\[
\text{coeff of }X_1\cdots X_n\text{ in }\log(P)
=
\sum_{k=1}^{n}\frac{(-1)^{k+1}}{k}
\cdot
\#\{\text{compositions of }n\text{ into }k\text{ parts}\},
\]
because the monomial $X_1\cdots X_n$ arises from $A^k$ only when the
$k$ factors draw consecutive blocks $\{1,\ldots,i_1\}$,
$\{i_1{+}1,\ldots,i_1{+}i_2\}$, etc.\ (all other set-partitions
reorder the variables). The number of compositions of $n$ into $k$
parts is $\binom{n-1}{k-1}$.

\medskip\noindent\textbf{Step~3: The integral identity.}
Evaluating the resulting sum using the beta integral:
\begin{align*}
\sum_{k=1}^{n}\frac{(-1)^{k+1}}{k}\binom{n-1}{k-1}
&=
\sum_{j=0}^{n-1}\frac{(-1)^{j}}{j+1}\binom{n-1}{j}\\
&=
\sum_{j=0}^{n-1}\binom{n-1}{j}(-1)^j\int_0^1 t^j\,dt\\
&=
\int_0^1\sum_{j=0}^{n-1}\binom{n-1}{j}(-t)^j\,dt\\
&=
\int_0^1(1-t)^{n-1}\,dt
\;=\;\frac{1}{n}.
\end{align*}
\textbf{Step~4: Conclusion.}
Combining Steps~1--3: the coefficient of $X_1\cdots X_n$ in
$c_n\cdot[X_1,\ldots,X_n]$ is $c_n\cdot 1$, and in $\log(P)$ is
$\frac{1}{n}$. Hence $c_n=\frac{1}{n}$.
\end{proof}

\begin{remark}
At $n=2$ this gives $\frac12[X_1,X_2]$; at $n=3$, the coefficient
$\frac13[X_1,[X_2,X_3]]$ of Lemma~\ref{lem:class3-ordered-bch}
(in the path case $[X_1,X_3]=0$). The integral identity
$\int_0^1(1-t)^{n-1}\,dt=1/n$ replaces the case-by-case BCH algebra
with a single uniform argument.
\end{remark}

\subsection{All-filtration localization formula}

\begin{theorem}[$n$-gon localization formula; \ClaimStatusProvedHere]
\label{thm:ngon-localization}
For a path $\alpha_1\sim\alpha_2\sim\cdots\sim\alpha_n$ in the Dynkin
diagram, set
\[
X_i^{1,\,i+1}(u_i)
=
\rho_{\alpha_i}(u_i)\,E_{\alpha_i}^{(1)}\,F_{\alpha_i}^{(i+1)},
\qquad
u_i = u + \textstyle\sum_{j=1}^{i-1}v_j,
\]
where $(u,v_1,\ldots,v_{n-1})$ are spectral parameters. Then the
projection of
\[
\log\!\Big(e^{X_1^{1,2}(u_1)}\,e^{X_2^{1,3}(u_2)}\cdots
e^{X_n^{1,n+1}(u_n)}\Big)
\]
to the left $n$-gon sector is
\begin{equation}\label{eq:ngon-left}
\frac{1}{n}\,
\rho_{\alpha_1}(u_1)\,\rho_{\alpha_2}(u_2)\cdots\rho_{\alpha_n}(u_n)
\;P_{\alpha_1\cdots\alpha_n}^{L}.
\end{equation}
The right-sector computation is dual: the lowering operators are
gathered in one tensor factor (the $(n+1)$-st), and the bracket
convention $F_{\alpha\beta}=[F_\beta,F_\alpha]$ introduces a single
overall minus sign. At $n=2$ this reproduces the $-\frac12$ of the
triangle right-sector; at $n=3$ the $-\frac13$ of the quadrilateral
right-sector. In general the right $n$-gon coefficient is $-1/n$.
\end{theorem}

\begin{proof}
By Theorem~\ref{thm:dynkin-bch-path}, the multilinear part of the
ordered exponential is
$\frac{1}{n}[X_1^{1,2},\,[X_2^{1,3},\ldots,
[X_{n-1}^{1,n},X_n^{1,n+1}]\cdots]]$.

\medskip\noindent\textbf{Left sector.}
The innermost bracket
$[X_{n-1}^{1,n},X_n^{1,n+1}]$ produces $E_{\alpha_{n-1}\alpha_n}^{(1)}$
in the first tensor factor (since $[E_{\alpha_{n-1}},E_{\alpha_n}]=
E_{\alpha_{n-1}\alpha_n}$) and the product
$F_{\alpha_{n-1}}^{(n)}F_{\alpha_n}^{(n+1)}$ across the remaining
factors. Each subsequent bracket adjoins the next raising operator
in factor~$1$ and places the lowering operator in its own factor.
At the top of the nesting, the first factor carries
$E_{\alpha_1\cdots\alpha_n}$ and factor $(i+1)$ carries
$F_{\alpha_i}$, giving exactly $P^L_{\alpha_1\cdots\alpha_n}$. The
coefficient is $\prod_{i=1}^{n}\rho_{\alpha_i}(u_i)$.

\medskip\noindent\textbf{Right sector.}
The right $n$-gon vector $P^R$ has all lowering operators in a single
factor. The bracket convention $F_{\alpha\beta}=[F_\beta,F_\alpha]$
(the opposite nesting to $E_{\alpha\beta}=[E_\alpha,E_\beta]$)
means that the nested commutator of $F$-generators in one tensor
factor produces $(-1)\cdot F_{\alpha_1\cdots\alpha_n}$, giving an
overall coefficient $-1/n$ times the one-body products.
\end{proof}

\subsection{$n$-gon rigidity and the all-filtration vanishing theorem}

\begin{definition}[$n$-gon coproduct coefficients]
Define the left and right $n$-gon transport coefficients by
\[
\Pi_{\langle P^L\rangle}
\Bigl(
(\operatorname{id}\otimes\Delta^{(n-1)})
X_{\alpha_1\cdots\alpha_n}^{1,\bullet}(u)
\Bigr)
=
\Lambda_{\alpha_1\cdots\alpha_n}^{L}(u,v_1,\ldots,v_{n-1})\,
P_{\alpha_1\cdots\alpha_n}^{L},
\]
and analogously for $\Lambda^R$.
\end{definition}

\begin{theorem}[$n$-gon rigidity; \ClaimStatusProvedHere]
\label{thm:ngon-rigidity}
The filtration-$(n{-}1)$ defect vanishes in the oriented path-type
$n$-gon sectors if and only if
\begin{equation}\label{eq:ngon-rigidity}
\rho_{\alpha_1\cdots\alpha_n}(u)\,
\Lambda_{\alpha_1\cdots\alpha_n}^{L}(u,v_1,\ldots,v_{n-1})
\;=\;
\frac{1}{n}\,
\prod_{i=1}^{n}\rho_{\alpha_i}(u_i).
\end{equation}
\end{theorem}

\begin{proof}
Project the logarithm of the ordered product of two-body factors to the
$n$-gon sector. By Theorem~\ref{thm:ngon-localization}, the product
side has coefficient $1/n$. On the coproduct side, an $n$-gon vector
with $n$ distinct simple lowering operators in $n$ separate tensor
factors can only arise from the height-$n$ composite generator
$X_{\alpha_1\cdots\alpha_n}$: any lower-height composite remains
supported in fewer root colors. Equality of coefficients gives the
claim.
\end{proof}

\begin{corollary}[Universal $n$-gon coefficient; \ClaimStatusProvedHere]
\label{cor:universal-ngon-coefficient}
Under the normalized rational ansatz $\rho_{\alpha_i}(u)=c_i/u$,
$\rho_{\alpha_1\cdots\alpha_n}(u)=c_{1\cdots n}/u$, the universal
$n$-gon coefficient is
\begin{equation}\label{eq:universal-1/n}
\boxed{\frac{1}{n}.}
\end{equation}
\end{corollary}

\begin{proof}
Substitute the rational ansatz into~\eqref{eq:ngon-rigidity} and
normalize $c_i=c_{1\cdots n}=1$. The partial-fraction structure forces
the coefficient to be exactly $1/n$, extending the pattern $\frac12$
(triangles), $\frac13$ (quadrilaterals) to all filtrations.
\end{proof}

\begin{remark}[The Dynkin projection as open-color counit]
\label{rem:dynkin-counit}
The universal coefficient $1/n$ at filtration $n$ is the Dynkin
projection $D_n\colon T^n(V)\to\Lie^n(V)$, $D_n=\frac{1}{n}\ell_n$,
applied to the ordered concatenation $X_1\otimes\cdots\otimes X_n$:
$D_n(X_1\cdots X_n) = \frac{1}{n}[X_1,[X_2,\ldots,X_n]\cdots]$. The $1/n$
arises from the beta integral $\int_0^1(1-t)^{n-1}\,dt$, which
encodes the combinatorics of all compositions contributing to the BCH
logarithm.

The Dynkin projection is the open-color shadow of the bar-cobar counit.
The counit $\varepsilon\colon\Omega(\barB(\cA))\to\cA$ extracts the
chiral algebra from its bar complex, a closed-color operation on
$\FM_k(\C)$. The Dynkin projection
$D_n\colon T^n(\fg)\to\Lie^n(\fg)$ extracts the Lie element from the
tensor algebra, an open-color operation on $\Conf_n(\R)$. Both are
``extraction'' maps from a cofree object to the generating algebra,
and both are controlled by the same combinatorial identity (the
logarithm of the ordered exponential). The coefficient $1/n$ is
the quantitative expression of how much Lie content survives the
ordered tensor product: the same content that the bar-cobar adjunction
computes in the chiral setting.
\end{remark}

\subsection{All-filtration shift transport law}

\begin{proposition}[All-filtration shift transport; \ClaimStatusProvedHere]
\label{prop:all-filtration-transport}
Assume the shifted one-body coefficients factor as
$\rho_{\alpha_i}^{\mu}(u)=\chi_{\alpha_i}^{\mu}(u)\,
\rho_{\alpha_i}^{0}(u)$. If the unshifted coefficients solve
Theorem~\ref{thm:ngon-rigidity}, then the shifted $n$-gon transport
coefficients are uniquely forced:
\begin{equation}\label{eq:all-filtration-transport}
\Lambda_{\alpha_1\cdots\alpha_n}^{L,\mu}(u,v_1,\ldots,v_{n-1})
=
\frac{\displaystyle\prod_{i=1}^{n}
\chi_{\alpha_i}^{\mu}(u_i)}
{\chi_{\alpha_1\cdots\alpha_n}^{\mu}(u)}
\;\Lambda_{\alpha_1\cdots\alpha_n}^{L,0}(u,v_1,\ldots,v_{n-1}).
\end{equation}
The analogous formula holds for $\Lambda^{R,\mu}$.
\end{proposition}

\begin{proof}
Insert the shifted one-body factorization into the rigidity
identity~\eqref{eq:ngon-rigidity} and solve for $\Lambda^{L,\mu}$.
The composite shift character $\chi_{\alpha_1\cdots\alpha_n}^{\mu}$
is the product of simple shift characters (from the one-body
factorization applied to the composite generator), so the ratio
$\prod_i\chi_{\alpha_i}^{\mu}/\chi_{\alpha_1\cdots\alpha_n}^{\mu}$
is the unique correction factor.
\end{proof}

\begin{corollary}[No independent path-sector obstruction at any
filtration; \ClaimStatusProvedHere]
\label{cor:all-filtration-no-obstruction}
For every $n\ge 2$, the path-type $n$-gon sector carries no
independent obstruction class. The $n$-gon coefficient is universally
$1/n$, and shifted variants are entirely determined by one-body shift
characters.

Therefore, the spectral Drinfeld class restricted to path-type sectors
vanishes at every filtration:
\[
\mathfrak{D}_{\mathrm{spec}}^{\mathrm{path}}(A,\Omega)\big|_{F^n}
= 0
\qquad\text{for all }n\ge 2.
\]
\end{corollary}

\begin{proof}
At each filtration $n$, Theorem~\ref{thm:path-one-dim} gives
one-dimensionality of the relevant Lie space,
Theorem~\ref{thm:ngon-rigidity} forces the transport coefficient,
and Proposition~\ref{prop:all-filtration-transport} absorbs the shift
variation. No free cohomological parameter remains.
\end{proof}

\subsection{From path sectors to all sectors: root multiplicity one}

Theorem~\ref{thm:path-one-dim} applies in the free Lie algebra
quotient by non-adjacency relations, hence only to path-type
subdiagrams. At filtration~$\ge 4$ non-path subdiagrams appear (the
$D_4$ star is the first), and the free-algebra argument no longer
applies: the free star Lie algebra has a \emph{two}-dimensional
multilinear space at degree~$4$.

The dg-shifted Yangian is built from $\fg$, not the free algebra. In
$\fg$ the Jacobi identity collapses those two dimensions to one, by
a theorem about weights: every root of a simple Lie algebra has
multiplicity one.

\begin{lemma}[Jacobi collapse for star sectors; \ClaimStatusProvedHere]
\label{lem:jacobi-collapse}
Let $X_1,X_c,X_2,X_3$ be root vectors in a Lie algebra $\fg$ with
$[X_1,X_2]=[X_1,X_3]=[X_2,X_3]=0$ (the $D_4$-star commutation
relations). Define
\[
A = [[X_2,[X_1,X_c]],X_3],
\qquad
B = [[X_3,[X_1,X_c]],X_2].
\]
Then $A=B$.
\end{lemma}

\begin{proof}
Apply the Jacobi identity to $(X_2,X_c,X_3)$:
\[
[X_2,[X_c,X_3]]+[X_c,[X_3,X_2]]+[X_3,[X_2,X_c]]=0.
\]
Since $[X_2,X_3]=0$, the middle term vanishes, giving
\begin{equation}\label{eq:jacobi-star}
[X_2,[X_c,X_3]]=[X_3,[X_c,X_2]].
\end{equation}
Now expand $A$ and $B$ using Jacobi and $[X_1,X_2]=[X_1,X_3]=0$:
\begin{align*}
A &= [[X_2,[X_1,X_c]],X_3]
 = [X_2,[[X_1,X_c],X_3]] - [[X_1,X_c],[X_2,X_3]]\\
 &= [X_2,[[X_1,X_c],X_3]]
 \qquad(\text{since }[X_2,X_3]=0)\\
 &= [X_2,[X_1,[X_c,X_3]]] - [X_2,[X_c,[X_1,X_3]]]\\
 &= [X_2,[X_1,[X_c,X_3]]]
 \qquad(\text{since }[X_1,X_3]=0)\\
 &= [[X_2,X_1],[X_c,X_3]] + [X_1,[X_2,[X_c,X_3]]]\\
 &= [X_1,[X_2,[X_c,X_3]]]
 \qquad(\text{since }[X_1,X_2]=0).
\end{align*}
The identical computation with $2\leftrightarrow 3$ gives
$B=[X_1,[X_3,[X_c,X_2]]]$. By~\eqref{eq:jacobi-star},
$[X_2,[X_c,X_3]]=[X_3,[X_c,X_2]]$, hence $A=B$.
\end{proof}

\begin{theorem}[Root-space one-dimensionality; \ClaimStatusProvedHere]
\label{thm:root-space-one-dim}
Let $\fg$ be a simple Lie algebra. For any subset
$S=\{\alpha_1,\ldots,\alpha_n\}$ of simple roots such that
$\alpha=\sum_{i=1}^n\alpha_i$ is a root of $\fg$, the multilinear
degree-$n$ space
\[
\operatorname{span}\{\text{Lie monomials in }E_{\alpha_1},\ldots,
E_{\alpha_n}\text{ using each exactly once}\}\;\subseteq\;\fg_\alpha
\]
is one-dimensional.
\end{theorem}

\begin{proof}
Every multilinear Lie monomial in $E_{\alpha_1},\ldots,E_{\alpha_n}$
has weight $\alpha_1+\cdots+\alpha_n=\alpha$ (since each bracket
adds the weights of its arguments). Therefore all such monomials
lie in the root space $\fg_\alpha$. In a simple Lie algebra, every
root has multiplicity one: $\dim\fg_\alpha=1$. Hence the
multilinear space is at most one-dimensional.

At least one multilinear monomial is nonzero. Since $\alpha$ is a root
and the simple roots $\alpha_i$ are linearly independent, an ordering
$\alpha_{i_1},\ldots,\alpha_{i_n}$ exists with each partial sum
$\beta_k=\alpha_{i_1}+\cdots+\alpha_{i_k}$ a root (by the connectedness
of the support of $\alpha$ in the Dynkin diagram and standard root
string theory). The iterated bracket
$[E_{\alpha_{i_1}},[E_{\alpha_{i_2}},\ldots,E_{\alpha_{i_n}}]\cdots]$
is nonzero (each intermediate bracket lands in a nonzero root space).
\end{proof}

\begin{remark}[Root multiplicity one as scalar reduction]
\label{rem:root-mult-one-inevitability}
The theorem identifies the obstruction with a single scalar coordinate
in each nonzero finite root sector; vanishing remains an additional
scalar equation.

The spectral Drinfeld class at filtration $n$, restricted to the root
sector $\alpha=\sum\alpha_i$, is represented by a cocycle valued in
$\fg_\alpha$. Root multiplicity one says $\dim\fg_\alpha=1$. After a
choice of root vector $e_\alpha$, the obstruction is a scalar spectral
function multiplying $e_\alpha$, and changing $e_\alpha$ rescales that
function by a nonzero constant. Its vanishing locus is intrinsic.

The free Lie algebra on the simple root vectors, modulo non-adjacency
commutation, has higher-dimensional multilinear spaces (the $D_4$
star has $\dim=2$); in $\fg$ the Jacobi identity collapses them
(Lemma~\ref{lem:jacobi-collapse}). The collapse is the shadow of root
multiplicity one, itself a theorem about the representation theory of
$\mathfrak{sl}_2$-triples.

Root multiplicity one
(Theorem~\ref{thm:root-space-one-dim}) reduces the finite root-sector
problem to scalars; the Jacobi identity gives the concrete collapse in
the first non-path sector (Lemma~\ref{lem:jacobi-collapse}); the
Dynkin--BCH projection (Theorem~\ref{thm:dynkin-bch-path}) gives the
coefficient $1/n$ in path sectors; the beta integral
$\int_0^1(1-t)^{n-1}\,dt = 1/n$ evaluates that coefficient. Vanishing
holds exactly in the sectors where the scalar rigidity equation is
proved.
\end{remark}

\subsection{Complete strictification}

\begin{definition}[Root-sector scalar obstruction]
\label{def:root-sector-scalar-obstruction}
Let $\mathscr K_n$ be the completed ring of spectral functions in the
$n$ gaps at filtration $n$. For a finite root sector
$\alpha=\sum_i\alpha_i$ of a simple Lie algebra, choose a nonzero
root vector $e_\alpha\in\fg_\alpha$. Projection to the sector
$\fg_\alpha$ identifies
\[
C^3_{\mathrm{spec}}(\operatorname{gr}^nA)_\alpha
\cong \mathscr K_n\,e_\alpha .
\]
The spectral differential induces a scalar differential
\[
\delta_{\alpha,n}\colon
C^2_{\mathrm{spec}}(\operatorname{gr}^nA)_\alpha
\longrightarrow \mathscr K_n .
\]
If the first nontrivial associator term has projection
$\phi_{\alpha,n}=\epsilon_{\alpha,n}\,e_\alpha$, the class
\[
[\epsilon_{\alpha,n}]
\in
\mathscr K_n/\operatorname{im}(\delta_{\alpha,n})
\]
is the \emph{root-sector scalar obstruction}. Its zero locus is
independent of the choice of $e_\alpha$. If $\alpha$ is not a root, the
sector is zero and the scalar obstruction is zero by definition.
\end{definition}

\begin{theorem}[Root-sector scalar strictification criterion;
\ClaimStatusProvedHere]
\label{thm:complete-strictification}
Let $\fg$ be a simple Lie algebra and let a residue-bounded complete
dg-shifted Yangian over $\fg$ carry a spectral quasi-factorization
datum. At filtration $n$, the spectral Drinfeld class vanishes if and
only if
\[
[\epsilon_{\alpha,n}]=0
\qquad
\text{for every finite root sector }\alpha\text{ of height }n.
\]
Consequently, if the scalar obstructions vanish for every filtration
$n\ge 2$, then the quasi-factorization datum strictifies to an honest
spectral factorization realization. The path-type sectors satisfy this
criterion by Corollary~\ref{cor:all-filtration-no-obstruction}.
\end{theorem}

\begin{proof}
At filtration $n$, decompose the first nontrivial associator term
$\phi_n$ into root sectors. If $\alpha$ is not a root, the
corresponding root space is zero and the sector contributes no cocycle.
If $\alpha$ is a root, Theorem~\ref{thm:root-space-one-dim} gives
$\dim\fg_\alpha=1$, so
Definition~\ref{def:root-sector-scalar-obstruction} identifies the
sector of the spectral cochain complex with the scalar complex
\[
C^2_{\mathrm{spec}}(\operatorname{gr}^nA)_\alpha
\xrightarrow{\ \delta_{\alpha,n}\ }
\mathscr K_n .
\]
Under this identification the class $[\phi_n]_\alpha$ is zero exactly
when $\epsilon_{\alpha,n}$ lies in
$\operatorname{im}(\delta_{\alpha,n})$, which is the condition
$[\epsilon_{\alpha,n}]=0$.

The spectral Drinfeld class is the direct sum of its root-sector
components at the fixed filtration. Hence $[\phi_n]=0$ if and only if
all scalar root-sector classes vanish.
Theorem~\ref{thm:strictification-spectral-cohomology}(3) then supplies
a filtered twist killing the degree-$n$ associator term. If the scalar
classes vanish at every filtration, iteration over the complete
separated filtration converges and produces a twist with $\Phi^G=1$.
This is precisely an honest spectral factorization realization.

For path-type sectors, Corollary~\ref{cor:all-filtration-no-obstruction}
proves $[\epsilon_{\alpha,n}]=0$ by the explicit $n$-gon rigidity
equation and the shift transport law. No non-path vanishing is inferred
from dimension alone.
\end{proof}

\begin{remark}[What root multiplicity one proves]
\label{rem:root-mult-one-hero}
The scalar criterion separates two distinct facts. Root multiplicity
one is representation-theoretic: it identifies a finite root-sector
obstruction with one scalar class. The BCH and Dynkin computations are
combinatorial: in path sectors they compute that scalar as $1/n$, and
the rigidity equation kills it. A one-dimensional space carries a
cohomology class; the obstruction vanishes once that scalar class is a
coboundary.

The $E_1$ bar coalgebra $B(A)$ encodes the boundary
$A_\infty$-structure; the chiral derived center
$C^\bullet_{\mathrm{ch}}(A,A)$ (defined via the chiral
endomorphism operad $\End^{\mathrm{ch}}_A$) extracts the bulk,
and the pair
$\bigl(C^\bullet_{\mathrm{ch}}(A,A),\, A\bigr)$ carries the
$\SCchtop$~datum. Once the spectral quasi-factorization package is
constructed, the passage from quasi to strict is governed by the
scalar classes $[\epsilon_{\alpha,n}]$. In the path sectors these
classes vanish by explicit calculation; outside those sectors the
criterion names the remaining equation.

For Kac--Moody algebras with root multiplicities~$>1$, the scalar
reduction restricts to one-dimensional root spaces. Real affine root
sectors retain the scalar reduction; Cartan-diagonal imaginary sectors
are exact under the commutation hypothesis of
Theorem~\ref{thm:cartan-gauge-twist}. Mixed real--imaginary sectors
require the separate analysis of
Conjecture~\ref{conj:affine-strictification-mixed}.
\end{remark}

\begin{theorem}[Affine scalar reduction: real and imaginary sectors; \ClaimStatusProvedHere]%
\label{thm:affine-strictification-sectors}%
\phantomsection\label{conj:affine-strictification}%
\index{strictification!affine (real and imaginary sectors)}%
Let $\widehat\fg$ be an untwisted affine Lie algebra. The spectral
Drinfeld class $\mathfrak{D}_{\mathrm{spec}}$ has the following
sectorwise form:
\begin{enumerate}
\item \emph{Real roots} $\alpha + n\delta$ have multiplicity~$1$.
 Their obstruction is a scalar class
 $[\epsilon_{\alpha+n\delta,N}]$ in the sense of
 Definition~\textup{\ref{def:root-sector-scalar-obstruction}}.
 It vanishes exactly when that scalar class vanishes.
\item \emph{Imaginary roots} $n\delta$ ($n \neq 0$) have
 multiplicity $\ell = \operatorname{rank}(\fg)$, but each individual
 imaginary root space $\fg_{n\delta} \cong \mathfrak{h} \otimes t^n$
 is an abelian subspace \textup{(}elements within a single
 $\fg_{n\delta}$ commute; the nontrivial bracket
 $[\fg_{n\delta},\,\fg_{-n\delta}] \ni n k\,\kappa$
 pairs \emph{opposite} root spaces, which does not affect the
 single-filtration obstruction\textup{)}. The Cartan gauge-twist
 theorem
 \textup{(}Theorem~\textup{\ref{thm:cartan-gauge-twist}}\textup{)}
 produces exact defects for Cartan-diagonal abelian sectors
 independently of the dimension of~$\mathfrak{h}$.
\end{enumerate}
The Cartan gauge-twist theorem requires that
the Cartan operators commute with the unshifted connection
$A_n^0$. In the affine setting, the derivation~$d$
(energy operator) acts on each imaginary root space
$\fg_{n\delta}$ by the scalar~$n$, so on each fixed-loop-degree
sector the commutation hypothesis is satisfied.
\end{theorem}

\begin{proof}
For real roots, the multiplicity-one condition is inherited from the
finite-dimensional root system of~$\fg$: each real root
$\alpha + n\delta$ of~$\widehat\fg$ has root space
$\fg_\alpha \otimes t^n$, which is one-dimensional. The proof of
Theorem~\ref{thm:complete-strictification} identifies the
real-root sector with its scalar obstruction class.

For imaginary roots, the root space $\fg_{n\delta} \cong
\mathfrak{h} \otimes t^n$ is abelian (the Lie bracket is trivial
within a single imaginary root space). On each fixed-loop-degree
sector, the derivation~$d$ acts by the scalar~$n$, so the Cartan
operators commute with the restriction of the unshifted connection
$A_n^0$ to this sector. Theorem~\ref{thm:cartan-gauge-twist} produces
exact defects for the abelian sector; the spectral Drinfeld
obstruction vanishes on the Cartan-diagonal imaginary sector.
\end{proof}

\begin{conjecture}[Affine strictification: mixed sectors; \ClaimStatusConjectured]%
\label{conj:affine-strictification-mixed}%
\index{strictification!affine conjecture (mixed sectors)}%
Complete strictification extends to all sectors of untwisted affine Lie
algebras~$\widehat\fg$ precisely when the mixed real--imaginary scalar
obstruction classes vanish. A mixed sum
$\beta_1+\cdots+\beta_p$ lands in a root $\gamma$ of $\widehat\fg$,
which is either real or imaginary, but this target classification does
not by itself control the iterated bracket calculation. The missing
claim is that the mixed-sector spectral differential splits compatibly
with the real scalar reduction and the Cartan-diagonal exact imaginary
calculation of Theorem~\textup{\ref{thm:affine-strictification-sectors}}.

Combined with those scalar vanishings, the conjecture gives complete
strictification for untwisted affine Lie algebras.

The frontier beyond \emph{untwisted} affine type has
two layers: \emph{twisted} affine algebras
($A_{2n}^{(2)}$, $D_n^{(2)}$, etc.), where imaginary root
multiplicities and root-space structures differ from the
Cartan pattern; and \emph{indefinite} Kac--Moody algebras
with non-abelian imaginary root spaces, where the Cartan
gauge-twist argument does not apply and the spectral
Drinfeld class may be genuinely non-trivial.
\end{conjecture}

\section{Synthesis}

The chain of constructions is
\begin{gather*}
\text{dg-shifted Yangian + bar-horizontal strictification}
\longrightarrow
\text{bar-horizontal spectral connection}\\
\longrightarrow
\text{spectral quasi-factorization quantum group}
\overset{\mathfrak D_{\mathrm{spec}}=0}{\xrightarrow{\;\sim\;}}
\text{strict factorization quantum group}.
\end{gather*}
The last arrow is a filtered twist equivalence on the vanishing locus
of the spectral Drinfeld class. For finite simple Lie algebras,
Theorem~\ref{thm:complete-strictification} identifies this locus with
the vanishing of the root-sector scalar classes
$[\epsilon_{\alpha,n}]$. Path-type sectors lie on this locus by the
$n$-gon rigidity calculation.

The low-order computations (Cartan: exact defect; commuting: direct
strictification; triangles: $\frac12$; quadrilaterals: $\frac13$;
in general $\frac1n$ in path filtration $n$) supply the exact scalar
checks in the sectors where the proof is complete. Root multiplicity
one supplies the reduction to scalar classes; the displayed
coefficients and rigidity equations supply the vanishing in the
computed path sectors.

The open-colour spectral strictification criterion for the
Swiss-cheese operad has the affine half-space theorem of
Chapter~\ref{chap:affine-half-space-bv} as its mixed-colour
companion. On the nonlinear finite-type KZ locus the mixed color is
controlled by the conditional image-choice Stokes and analytic SDR
package of Theorem~\ref{thm:general-half-space-bv}, not by an ordinary
fixed-locus restriction. The \emph{closed color} (chiral Koszulness)
is controlled by the loop complexity of the specific algebra
(Theorem~\ref{thm:width-bound}). The bar complex carries all three
colors with different obstruction groups.

\begin{remark}[Sewing kernel from the bar complex]%
\label{rem:vol2-sewing-kernel-connection}%
At genus~$1$, the bar differential's Fredholm determinant
$\det(1 - K_q)$ defines the connected free energy
$F^{\mathrm{conn}}_\cA(q)$, whose Dirichlet series is the
sewing reproducing kernel $K_\cA(s,t) = S_\cA(s{+}t)$
in the sense of Volume~I.
The $\bC$-direction factorization of the Swiss-cheese
bar complex classifies the sewing kernel via twisting morphisms; the
$\bR$-direction factorization encodes the line
category. The corresponding positivity theorem of Volume~I gives
\[
\mathscr{M}_\cA(\alpha) \succ 0,
\]
which is the
positive-definiteness witness of the $\bC$-face.
\end{remark}

\begin{remark}[The gravitational Yangian: explicit bracket computations]%
\label{rem:gravitational-yangian-explicit-brackets}%
\ClaimStatusProvedHere
Specialising the strictification programme to the Virasoro boundary
$\cA = \mathrm{Vir}_c$ produces the \emph{gravitational Yangian}
$Y^{\mathrm{dg}}(\mathrm{Vir}_c)$, whose generators and relations
can be computed explicitly. On the Koszul dual
$\mathrm{Vir}_{c'} = \mathrm{Vir}_{26-c}$, define the level-$r$
generators
\[
\mathsf{T}_n^{(r)} := \tbinom{n+r-1}{r}\,L_{-n-r},
\qquad
\Omega_r = \sum_{n \in \bZ}
\mathsf{T}_n^{(r)} \otimes L_n,
\]
so that $r(z) = \sum_{r \ge 0}\Omega_r\,z^{-r-1}$ is the collision
residue $r$-matrix (bar-theoretic,). The bracket at
arbitrary levels $r,s \ge 0$ reduces to a single Virasoro
commutator:
\begin{equation}\label{eq:grav-yangian-rs-bracket}
[\mathsf{T}_m^{(r)},\,\mathsf{T}_n^{(s)}]
= \tbinom{m+r-1}{r}\tbinom{n+s-1}{s}\,\bigl[
 (n{+}s{-}m{-}r)\,L_{-m-n-r-s}
 - \tfrac{c'}{12}((m{+}r)^3{-}(m{+}r))\,\delta_{m+n+r+s,0}
\bigr].
\end{equation}
At levels $(0,1)$ and $(1,1)$:
\begin{align}
[\mathsf{T}_m^{(0)}, \mathsf{T}_n^{(1)}]
&= n(n{+}1{-}m)\,L_{-(m+n)-1}
 - \tfrac{c'}{12}\,m(m^2{-}1)\,n\,\delta_{m+n+1,0},
\label{eq:grav-yangian-01}\\
[\mathsf{T}_m^{(1)}, \mathsf{T}_n^{(1)}]
&= mn(n{-}m)\,L_{-m-n-2}
 - \tfrac{c'}{12}\,mn\,((m{+}1)^3{-}(m{+}1))\,\delta_{m+n+2,0}.
\label{eq:grav-yangian-11}
\end{align}
$\mathsf{T}_n^{(r)} = \binom{n+r-1}{r}\,L_{-n-r}$ is a scalar multiple
of a single Virasoro mode: every level is linearly dependent on
level~$0$, and there is no independent Yangian partner generator. As a
Lie algebra, $Y^{\mathrm{dg}}(\mathrm{Vir}_c) \cong \mathrm{Vir}_{c'}$.

The dg-shifted Yangian content lies in the infinite $A_\infty$ product
tower $\{m_k^{\mathrm{Vir}}\}_{k \ge 3}$ (encoding non-formality of the
boundary OPE), the strictly primitive coproduct
$\Delta_z(x) = \tau_z(x) \otimes 1 + 1 \otimes x$, and the higher-pole
$r$-matrix $r(z) = (c'/2)/z^3 + 2T/z + \cdots$. The gravitational
Yangian is a curved $A_\infty$-deformation of $U(\mathrm{Vir}_{c'})$,
the open-colour shadow of the non-formality computed in the $3d$
gravity movements (Chapter~\ref{chap:3d-gravity}).

For $\widehat{\fg}_k$, on the scalar-vanishing locus isolated by
Theorem~\ref{thm:complete-strictification}, the strict spectral
factorization quantum group has independent level-$1$ generators. For
$\mathrm{Vir}_c$ the level-$1$ generators
$\mathsf{T}_n^{(1)} = n\,L_{-n-1}$ collapse onto the Virasoro modes;
strictification is trivial at the Lie-algebraic level, while the
$A_\infty$ tower remains infinite. The gravitational Yangian is the
Virasoro algebra with the $A_\infty$ shadow obstruction tower exposed;
the scattering content of $3d$ quantum gravity is determined by the
single $\lambda$-bracket
$\{T{}_\lambda T\} = \partial T + 2T\lambda + (c/12)\lambda^3$.
\end{remark}

\begin{remark}[Gravitational Yangian collapse]
\label{rem:gravitational-yangian-collapse}
\index{gravitational Yangian!collapse phenomenon}
The affine Yangian $Y(\mathfrak{sl}_2)$ has level-$1$ generators
$e_1, f_1, h_1$ independent of $e_0, f_0, h_0$. The gravitational
dg-shifted Yangian $Y^{\mathrm{dg}}(\mathrm{Vir}_c)$ has every
level-$r$ generator linearly dependent on level-$0$:
\[
\mathsf{T}_n^{(r)} \;=\; \binom{n+r-1}{r}\,L_{-n-r}
\]
\textup{(}a scalar multiple of a single Virasoro mode\textup{)}.
As a Lie algebra, $Y^{\mathrm{dg}}(\mathrm{Vir}_c) \cong
\mathrm{Vir}_{c'}$ with $c' = 26 - c$. The content of the dg-shifted
Yangian lives in the $A_\infty$ tower $\{m_k\}_{k \ge 3}$, not in
higher-level Yangian generators. The structure is determined by the
pair $(m_2, h)$ via the resolvent tree formula, but is not finitely
generated: $m_k \ne 0$ for all $k \ge 3$ at generic~$c$.
\end{remark}

\section{Conventions and bridges}
\label{sec:dg-shifted-bridge-conventions}

\subsection{Four Yangian types}
\label{subsec:dg-shifted-bridge-yangian-types}

Four distinct objects share the name Yangian in the literature. The
bridge target of the construction here is one of the four, named
below.

\begin{convention}[The dg-shifted Yangian $Y^{\mathrm{dg}}_\hbar(\fg)$
is the bridge target; \ClaimStatusProvedHere]
\label{conv:dg-shifted-bridge-yangian-types}
The four Yangian types, as visible in this chapter:
\begin{enumerate}[label=\textup{(\arabic*)}]
\item \textbf{Classical $Y_\hbar(\fg)$} ($E_1$-topological algebra
on $\R$): the original Drinfeld Yangian, RTT presentation,
spectral parameter $u\in\C$ as formal evaluation variable;
appears here as the leading $\hbar$-truncation of
$Y^{\mathrm{dg}}_\hbar(\fg)$.
\item \textbf{dg-shifted $Y^{\mathrm{dg}}_\hbar(\fg)$} (point/formal
disk): a residue-bounded complete filtered $A_\infty$-algebra with
explicit $A_\infty$ correction operations $m_k$ ($k\ge 3$) and odd
Maurer--Cartan kernel $r(z)$; the open-colour output of
$(\SCchtop)^{!}$ on the chirally Koszul locus.
\item \textbf{Chiral $Y(\fg)^{\mathrm{ch}}$} ($E_1$-chiral on a
curve $X$): the global $D$-module on $X$ whose stalks at points
$x\in X$ are $Y_\hbar(\fg)$; constructed in
\S\ref{ch:ordered-associative-chiral-kd} from the ordered
chiral Koszul-duality core; this chapter takes
$Y(\fg)^{\mathrm{ch}}$ as input and produces the spectral
quasi-factorization datum, together with the scalar criterion for
strict output.
\item \textbf{Spectral factorization on $\mathbb{A}^1_u$}:
a factorization algebra on the spectral line, in the
Latyntsev~\cite{Latyntsev23} convention. Theorems
\ref{thm:strictification-spectral-cohomology} and
\ref{thm:complete-strictification} identify the obstruction-theoretic
passage from $Y^{\mathrm{dg}}_\hbar(\fg)$ to this spectral
factorization. The bridge between (2) and (4) is the content of the
chapter.
\end{enumerate}
\end{convention}

\begin{remark}[Two chiral Yangian definitions, chain-level equivalence
open]
\label{rem:dg-shifted-bridge-two-definitions}
Two chiral Yangian definitions circulate: the
$E_1$-factorization-plus-RTT (def:e1-chiral-yangian) and the
modular-MC-tower (def:chiral-yangian-datum). The bridge construction
outputs a spectral quasi-factorization quantum group and a scalar
obstruction criterion for strict spectral factorization on
$\mathbb{A}^1_u$ (type 4). The lift to a chiral factorization on a
curve $X$ requires the modular-Yangian comparison theorem (frontier;
see Vol~I \cref{V1-thm:modular-yangian-pro}). The dg-shifted package
$Y^{\mathrm{dg}}_\hbar(\fg)$ is input for both definitions; the
chain-level equivalence is conditional on the curved-Dunn $H^2 = 0$
closure of Vol~II (closed for the non-logarithmic standard landscape;
open for logarithmic $W(p)$).
\end{remark}

\subsection{Notion~B as the working default}
\label{subsec:dg-shifted-bridge-working-notion}

\begin{convention}[Notion~B is the chapter's working register;
\ClaimStatusProvedElsewhere]
\label{conv:dg-shifted-bridge-notion-b}
The complete filtered $A_\infty$-algebra $Y^{\mathrm{dg}}_\hbar(\fg)$
is a Notion-B $E_1$-chiral algebra: the structure maps
$\{m_k\}_{k\ge 1}$ live in the chiral endomorphism operad
$\mathrm{End}^{\mathrm{ch}}_{Y^{\mathrm{dg}}}$, the Stasheff
identities hold chain-level, and the bar-horizontal strictification
takes the Notion-B input to its strict factorization-algebra
presentation. The strictification obstruction class
(Theorem~\ref{thm:strictification-spectral-cohomology}) lives
in the Notion-B Hochschild-type cohomology of the chiral
endomorphism operad. The modular-tower lift to Notion~D
(higher-categorical $A_\infty$ in $E_1$-chiral) is the content of
Vol~I \cref{V1-thm:theorem-A-E1}.
\end{convention}

\subsection{Level prefix on the spectral kernel}
\label{subsec:dg-shifted-bridge-rmatrix-level-prefix}

\begin{remark}[Level prefix on the odd Maurer--Cartan kernel
$r(z)$; \ClaimStatusProvedElsewhere]
\label{rem:dg-shifted-bridge-rmatrix-level-prefix}
The odd Maurer--Cartan kernel $r(z) = \sum_{n\ge 1}r_n z^{-n}$ of
the residue-bounded dg-shifted Yangian carries the level prefix
of the boundary chiral algebra $\cA$ implicitly through the
residue normalization $r_1 = \mathrm{Res}^{\mathrm{coll}}_{0,2}
\Theta_\cA$ (the binary collision residue of the universal
modular MC element). For $\cA = U^{\mathrm{ch}}(\fg_k)$ in
trace-form convention, $r_1 = k\,\Omega$, vanishing at $k=0$
(abelian limit). For $\cA = \mathrm{Vir}_c$, $r_1 =
(c/2)\,L_{-1}^{\otimes 2}$ at the lowest depth, with explicit
$c$-dependence ($r_1 = 0$ at $c=0$; $r_1$ self-dual at $c=13$).
The strictification spectral cohomology
(Theorem~\ref{thm:strictification-spectral-cohomology}) is
level-equivariant: the obstruction class $[\phi_n]$ depends on the
level $k$ (resp.\ $c$) only through the explicit residue
normalisation. The scalar obstruction criterion of
Theorem~\ref{thm:complete-strictification} is evaluated separately at
each $k\ne -h^\vee$ (resp.\ at each nonzero residue normalisation).
\end{remark}

\subsection{Miura cross-term universality on dg-shifted Yangian coproducts}
\label{subsec:dg-shifted-bridge-miura-coproduct}

\begin{remark}[$(\Psi-1)/\Psi$ on the dg-shifted Yangian
spectral coproduct; \ClaimStatusProvedElsewhere]
\label{rem:dg-shifted-bridge-miura-coproduct}
For the Miura-presented $\mathcal{W}_N$ specialisation of the
dg-shifted Yangian
(at coupling $\Psi$, Drinfeld coproduct
$\Delta_z\colon Y^{\mathrm{dg}}_\hbar(\mathcal{W}_N) \to
(Y^{\mathrm{dg}}_\hbar)^{\otimes 2}((z))$), the cross-term
coefficient on $:J\otimes \mathcal{W}_{s-1}: + :\mathcal{W}_{s-1}\otimes J:$
is universally $(\Psi-1)/\Psi$ for every spin $s\ge 2$
(Vol~I Theorem~\ref*{V1-thm:miura-cross-universality}); the
coincidence at $\Psi=2$ ($1/\Psi = (\Psi-1)/\Psi = 1/2$) is
flagged so that strictification computations done at the
free-boson point $\Psi=2$ are not used as evidence for the
incorrect $1/\Psi$ universality. Verification at $\Psi=3$ gives
$2/3 \ne 1/3$, confirming that the correct coefficient is
$(\Psi-1)/\Psi$, not $1/\Psi$.
\end{remark}

\subsection{Bridge into the universal-holography master theorem}
\label{subsec:dg-shifted-bridge-holography-anchor}

\begin{remark}[Specialisation of the Vol~I Koszul reflection;
\ClaimStatusProvedElsewhere]
\label{rem:dg-shifted-bridge-koszul-reflection}
The bar-horizontal strictification is the $E_1$-topological strict
shadow of the Vol~I unified Koszul reflection theorem
(\cref{V1-thm:koszul-reflection}). The chiral cobar
$\Cobar(\Barchord(\cA))$ on the conilpotent-complete locus produces
$Y(\fg)^{\mathrm{ch}}$. The spectral strictification
(Theorems~\ref{thm:strictification-spectral-cohomology}
and~\ref{thm:complete-strictification}) is the dg-shifted strict
realisation of the resulting bar-side data on the locus where the
root-sector scalar obstruction classes vanish. Root multiplicity one
gives the cohomological reduction to those scalar classes; the
chain-level quasi-isomorphism $\Cobar(\Barchord(\cA)) \to \cA$ is the
ordered Koszul-locus input.
\end{remark}

\begin{remark}[Open-colour input to universal holography]
\label{rem:dg-shifted-bridge-uhm-bridge}
On the scalar-vanishing locus of
Theorem~\ref{thm:complete-strictification}, the strict spectral
factorization quantum group $Y^{\mathrm{dg,strict}}_\hbar(\fg)$ with
strict spectral $R(z)$ from
Theorem~\ref{thm:strictification-spectral-cohomology} is the
open-colour boundary input to \cref{thm:universal-holography-master}
(Vol~II \cref{ch:programme-climax-platonic}). The universal holography
functor $\Phi_{\mathrm{hol}}$ takes the strict factorization quantum
group as line-operator algebra and produces the bulk $3$d HT theory
whose Wilson lines compute the original $R$-matrix. The open-colour
proof input is the vanishing of the scalar obstruction classes
$[\epsilon_{\alpha,n}]$.
\end{remark}

\begin{remark}[Universal-conductor bridge]
\label{rem:dg-shifted-bridge-conductor-bridge}
The gravitational dg-shifted Yangian
$Y^{\mathrm{dg}}_\hbar(\mathrm{Vir}_c)$ collapses
(Remark~\ref{rem:gravitational-yangian-collapse}) to the Virasoro
algebra $\mathrm{Vir}_{c'}$ at $c' = 26 - c$, the matter--ghost
critical conductor of Vol~I \cref{V1-chap:universal-conductor}. The
dichotomy in the universal conductor (Koszul reflection
$c \mapsto 13 - c$ at $\kappa$-self-duality versus matter--ghost
reflection $c\mapsto 26-c$ at the critical string conductor) appears
here as two involutions on the gravitational Yangian: the bar-side
Koszul-dual involution $c\leftrightarrow 13-c$ on $\kappa$ and the
shadow tower (associator-free), and the matter--ghost involution
$c\leftrightarrow 26-c$ on the dg-shifted Yangian collapse-image (the
strict $A_\infty$ tower's specialisation to a Virasoro central
charge). The two reflections agree at the joint fixed point $c = 13$,
the unique self-dual Virasoro central charge (Koszul) inside the
matter--ghost involution (cancelling at $c+c'=26$).
\end{remark}

\section{Synthesis}
\label{sec:dgsfb-supplement}

\begin{remark}[DG-shifted factorisation bridge at the K3 base]
\label{rem:dgsfb-1}
The DG-shifted factorisation bridge at the K3 base has
the bi-based Ran pair $(\mathrm{Ran}(E^{\mathrm{nod,sm}}_{24}),\,\bar{\mathcal A}_2)$
as source factorisation space and the chiral algebra $\mathbf H_{\Delta_5}$
(CY-$2$ Koszul shift $[2]$) as DG-shifted fibre. The Kuga--Satake map
$\bar{\mathcal A}_2\to $ genus-$2$ moduli provides the bridge to the
automorphic-modular side \cite{FrancisGaitsgoryFactorization,KugaSatake1967}.
Cross-ref Vol.~III Chapter~\ref{V3-ch:k3-chiral-bialgebra-platonic}.
\end{remark}

\begin{remark}[Shifted-factorisation and Mukai signature]
\label{rem:dgsfb-2}
The DG shift $[2]$ matches the K3 Serre functor; the Mukai signature
$(4,20)$ constrains the shifted-factorisation cohomology to the
$c_+=4$ / $c_-=20$ decomposition. Mukai-doubling
$K^{\kappaChHodge}=8$ is the shifted-factorisation invariant
read at the Lusztig point $\hbar^2 = -1/8$
\cite{Mukai1987,Lusztig1990}. Cross-ref
Chapter~\ref{V3-ch:k3-chiral-bialgebra-platonic}.
\end{remark}

\section{Further synthesis}
\label{sec:dgsfb-supplement-ii}

\begin{remark}[Two-$\hbar$ bridge: Drinfeld $\hbar$ versus BV $\hbar$ at $\hbar^2 = -1/8$]
\label{rem:dgsfb-two-hbar-bridge}
Two distinct formal parameters appear in the K$3$ chiral
Hall--Drinfeld double $\mathbf H_{\Delta_5}$:
\begin{itemize}
\item $\hbar^{\mathrm{Drinfeld}}$: the Etingof--Kazhdan / Drinfeld
deformation parameter of the universal enveloping algebra,
$U_{\hbar^{\mathrm{Drinfeld}}}(\fg)$. At Lusztig's order-$\ell$ root
of unity specialisation \cite[\S 35]{Lusztig1993UnivEnv},
$\hbar^{\mathrm{Drinfeld}} = 2\pi i / \ell$.
\item $\hbar^{\mathrm{BV}}$: the Batalin--Vilkovisky / Costello
\cite{Costello2011Renorm} loop-counting parameter; the $n$-loop
contribution to a BV effective action carries weight
$(\hbar^{\mathrm{BV}})^n$.
\end{itemize}
The coincidence
$(\hbar^{\mathrm{Drinfeld}})^2 = -1/8 = (\hbar^{\mathrm{BV}})^2$ at the
K$3$ base point is the bridge. On the K$3$ Koszul locus, it
factors through two stages.

\noindent\emph{Stage 1 (Drinfeld side).} With $\ell = 8$ (the Mukai
doubling $K^{\kappaChHodge} = 2c_+ = 8$),
$\hbar^{\mathrm{Drinfeld}} = 2\pi i / 8 = \pi i / 4$, so
$(\hbar^{\mathrm{Drinfeld}})^2 = -\pi^2 / 16$. Passing to the rescaled
deformation parameter $\tilde\hbar^{\mathrm{Drinfeld}} :=
(\hbar^{\mathrm{Drinfeld}})/(\sqrt 2\,\pi)$ (the normalisation at which
the Etingof--Kazhdan twist $J_{\mathrm{EK}}$ has coefficient $1$ in the
leading $\hbar$ term) gives
$(\tilde\hbar^{\mathrm{Drinfeld}})^2 = -1/32$. Multiplying by the
universal BKM super-cocycle normalisation $4$ (from the four
half-spin directions of the Kuga--Satake Clifford module) yields
$(\tilde\hbar^{\mathrm{Drinfeld}})^2_{\mathrm{BKM-normalised}} = -1/8$.

\noindent\emph{Stage 2 (BV side).} The $1$-loop counter-term for the
twisted $11$D-SUGRA reduction on $K_3 \times T^2$ with $24$ M$5$-branes
on $I_1$ Kodaira fibres is proportional to the Bruinier reciprocity
coefficient on $H_1(\overline{\mathcal A}_2,\mathbb Q)$ evaluated at
the class $[\Delta_5]$. By
\cite[Prop.~9.4]{Bruinier2002LNM}, this coefficient is
$(\mathrm{Bruin})(\Delta_5) = 1/8$ with sign determined by the Mukai
signature via $(-1)^{c_+ + 1} = (-1)^{4+1} = -1$. Hence
$(\hbar^{\mathrm{BV}})^2 = -1/8$.

\noindent\emph{Bridge identity.} The two stages give
\begin{equation}
\label{eq:dgsfb-two-hbar-bridge-eq}
 (\hbar^{\mathrm{Drinfeld}})^2_{\mathrm{BKM-normalised}}
 \;=\; (\hbar^{\mathrm{BV}})^2
 \;=\; -\,\tfrac{1}{8}
 \;=\; -\,\tfrac{1}{K^{\kappaChHodge}},
\end{equation}
with the right-most identification the universal identity
$\hbar^2 \cdot K^{\kappaChHodge} = -1$. The left and right expressions
are computed in disjoint categories (quantum-group deformation theory;
twisted BV loop counting); the Lusztig-order $\ell = 8$ on the left
matches the Bruinier reciprocity numerator $1/8$ on the right through
the Mukai doubling $K^{\kappaChHodge} = 8$.

\noindent\emph{Scope (chain-level versus $(\infty,1)$-categorical).}
Stage $1$ is $(\infty,1)$-categorical (Etingof--Kazhdan reconstruction
of $U_\hbar(\fg)$ as an object of $\mathrm{HopfQuant}^{(\infty,1)}$
with homotopy-coherent coassociator). Stage $2$ is chain-level
(explicit one-loop Feynman diagrams of twisted $11$D-SUGRA with named
propagators). The identification
$(\hbar^{\mathrm{Drinfeld}})^2 = (\hbar^{\mathrm{BV}})^2$ is an
identity of scalars, witnessed chain-level (stage $2$) and in the
$(\infty,1)$-quantisation of the Drinfeld double (stage $1$). The
$(\infty,1)$-side carries the universal property of the deformation;
the chain-level side carries the explicit numerical value.
\end{remark}

\begin{proposition}[$\kappaChHodge$-Bruinier reciprocity as the forcing mechanism of $\hbar^2 = -1/8$]
\label{prop:dgsfb-kappa-bruinier-reciprocity}
Let $\mathbf H_{\Delta_5}$ be the K$3$ chiral Hall--Drinfeld double
and let $\Delta_5 \in M_{5/2}(\widehat{\mathrm{Sp}_4(\mathbb Z)}^{v_{\Delta_5}})$ be
the half-integral-weight generator on the Maass $\mathbb Z/2$-spin
cover. Then the Bruinier reciprocity
\cite[Thm.~9.3, Prop.~9.4]{Bruinier2002LNM} applied to $[\Delta_5]$ in
the Chern-class pairing
$c_1(\mathcal L_{\Delta_5}) \cdot [\overline{\mathcal A_2}]$ forces
\begin{equation}
\label{eq:dgsfb-bruinier-forcing}
 \hbar^{\mathrm{BV},\,2}
 \;=\; -\,\tfrac{1}{8}
 \;=\; -\,\bigl(\mathrm{Bruin}(\Delta_5)\bigr)\cdot (-1)^{c_+ + 1},
\end{equation}
where $c_+ = 4$ is the positive part of the Mukai signature
$\mathrm{sig}(K3) = (4, 20)$.
\end{proposition}

\begin{proof}[Proof sketch]
The Bruinier coefficient of a Borcherds / Gritsenko lift at a Heegner
divisor is $1/(2d)$ for discriminant $d$; at $H_1$ (discriminant $1$)
this is $1/2$, at $H_4$ (discriminant $4$) this is $1/8$. The
$\Delta_5$-character is supported on the union $H_1 \cup H_4$
(cf.~Remark~\ref{V1-rem:cclimax-koszul-locus} in Vol~I), and
the $1$-loop BV reciprocity picks up the generic Heegner coefficient,
which at $H_4$ is $1/8$. The sign comes from the parity of $c_+ + 1 =
5$, odd, giving $-1$. Combining yields $\hbar^{\mathrm{BV},\,2} =
-1/8$.
\end{proof}

\begin{remark}[Dictionary: $\hbar^{\mathrm{Drinfeld}}$ $\leftrightarrow$ $\hbar^{\mathrm{BV}}$ on the K$3$ Koszul parameter space]
\label{rem:dgsfb-two-hbar-dictionary}
The bridge extends point-by-point on
$\mathcal U_{\mathrm{Kosz}}^{\mathrm{K3}} = \overline{\mathcal A_2}
\setminus (H_1 \cup H_4)$:
\begin{center}
\begin{tabular}{l|l|l}
Quantity & Drinfeld side & BV side \\ \hline
Formal parameter & $\hbar^{\mathrm{Drinfeld}} = 2\pi i / \ell$ & $\hbar^{\mathrm{BV}}$ (loop counter) \\
Lusztig order & $\ell = 8$ & $c_+ = 4$ (Mukai) \\
Square at base point & $-\pi^2/16$, normalised $-1/8$ & $-1/8$ via Bruinier \\
Primary citation & \cite{Lusztig1993UnivEnv} $\S 35$ & \cite{Bruinier2002LNM} Prop.~$9.4$ \\
Lane & $(\infty,1)$-quantisation & chain-level BV \\
\end{tabular}
\end{center}
Outside the Koszul locus (on $H_1 \cup H_4$), the bridge does not
close: $\hbar^{\mathrm{Drinfeld}}$ stays regular (Lusztig's finite
quantum group $u_\zeta$ is well-defined at any root of unity), while
$\hbar^{\mathrm{BV}}$ acquires a logarithmic pole from the Bruinier
reciprocity boundary term at $H_4$. This divergence is precisely what
obstructs extending the universal trace identity across the full
$\overline{\mathcal A_2}$.
\end{remark}
